\documentclass{article}

\usepackage[preprint]{tmlr}
\usepackage{amsmath}
\usepackage{amssymb}
\usepackage{mathtools}
\usepackage{longtable,array}
\usepackage{booktabs}
\usepackage{siunitx}
\sisetup{round-mode=places, round-precision=2, table-number-alignment=center, group-separator={}}
\usepackage{newunicodechar}
\newunicodechar{Σ}{\ensuremath{\Sigma}}
\newunicodechar{σ}{\ensuremath{\sigma}}
\newunicodechar{δ}{\ensuremath{\delta}}
\newunicodechar{α}{\ensuremath{\alpha}}
\newunicodechar{Ω}{\ensuremath{\Omega}}
\newunicodechar{ε}{\ensuremath{\varepsilon}}
\newunicodechar{γ}{\ensuremath{\gamma}}
\newunicodechar{ρ}{\ensuremath{\rho}}
\newunicodechar{λ}{\ensuremath{\lambda}}
\newunicodechar{η}{\ensuremath{\eta}}
\newunicodechar{μ}{\ensuremath{\mu}}
\newunicodechar{κ}{\ensuremath{\kappa}}
\newunicodechar{ν}{\ensuremath{\nu}}
\newunicodechar{φ}{\ensuremath{\phi}}
\newunicodechar{Ψ}{\ensuremath{\Psi}}
\newunicodechar{Ξ}{\ensuremath{\Xi}}
\newunicodechar{Θ}{\ensuremath{\Theta}}

\usepackage{graphicx}
\usepackage{tikz}
\usepackage{pgfplots}
\pgfplotsset{compat=1.18}
\usetikzlibrary{patterns, positioning, arrows.meta, shapes.geometric, backgrounds, fit}

\usepackage{algorithm}
\usepackage{algpseudocode}
\usepackage{enumitem}
\usepackage{amsthm}
\usepackage{microtype}
\usepackage{hyperref}
\hypersetup{
  pdftitle={Critical Compositional Pressure: A Phase-Boundary Framework for Compositional Representation Formation in Neural Networks},
  pdfsubject={Machine Learning, Dynamical Systems, Representation Learning},
  pdfkeywords={compositional generalization, continuous gradient flow, phase transition, shortcut learning, transcritical bifurcation}
}
\usepackage{xr}
\externaldocument{manuscript_journal}
\usepackage[capitalize,noabbrev]{cleveref}
\theoremstyle{plain}
\newtheorem{theorem}{Theorem}
\newtheorem{proposition}{Proposition}
\newtheorem{lemma}{Lemma}
\newtheorem{corollary}{Corollary}

\theoremstyle{definition}
\newtheorem{definition}{Definition}
\newtheorem{conjecture}{Conjecture}
\newtheorem{assumption}{Assumption}
\newtheorem{remark}{Remark}

\newcommand{\R}{\mathbb{R}}
\DeclareMathOperator{\Var}{Var}
\DeclareMathOperator{\Corr}{Corr}
\DeclareMathOperator{\argmax}{arg\,max}
\DeclareMathOperator{\CosSim}{CosSim}
\DeclareMathOperator{\CKA}{CKA}
\DeclareMathOperator{\GCA}{GCA}
\DeclareMathOperator{\Tr}{Tr}
\DeclareMathOperator{\diag}{diag}

% Okabe-Ito Colorblind-Safe Palette
\definecolor{okabe-black}{RGB}{0,0,0}
\definecolor{okabe-orange}{RGB}{230,159,0}
\definecolor{okabe-skyblue}{RGB}{86,180,233}
\definecolor{okabe-bluishgreen}{RGB}{0,158,115}
\definecolor{okabe-yellow}{RGB}{240,228,66}
\definecolor{okabe-blue}{RGB}{0,114,178}
\definecolor{okabe-vermillion}{RGB}{213,94,0}
\definecolor{okabe-purple}{RGB}{204,121,167}

\title{Supplementary Materials: Critical Compositional Pressure: A Phase-Boundary Framework for Compositional Representation Formation in Neural Networks\thanks{All mathematical proofs, analytical derivations, empirical protocols, raw seed tables, and econometric diagnostic tables are self-contained in this supplementary document. Open-source code and data repository: \url{https://github.com/basyirin-dev/sigma-model} (Release \texttt{v2.0-paper02}, Commit \texttt{540c992}). Complete reproducibility bundle: \texttt{supplementary\_materials.zip}.}}

\author{\name Basyirin Amsyar Basri \email basyirinbasri@gmail.com \\
      \addr Independent Researcher}

\begin{document}

\maketitle

\appendix

\section{Mathematical Notation \& Taxonomy Reference}
\label{sec:appendix_notation}

Table~\ref{tab:notation_taxonomy} provides a comprehensive taxonomy of mathematical state variables, coupling constants, order parameters, and diagnostic metrics utilized throughout the manuscript.

\begin{table}[h]
\centering
\caption{\textbf{Comprehensive Notation and Variable Taxonomy Reference.}}
\label{tab:notation_taxonomy}
\vskip 0.1in
\begin{small}
\begin{tabular}{lll}
\toprule
Symbol & Description & Domain / Units \\
\midrule
\multicolumn{3}{l}{\textbf{State Variables \& Dynamical Coordinates}} \\
$u(t)$ & Effective shortcut representation coordinate & $u \in \mathbb{R}_{\ge 0}$ \\
$v(t)$ & Coherent schema representation coordinate & $v \in \mathbb{R}_{\ge 0}$ \\
$\theta_S$ & Nominal empirical shortcut target & $\theta_S \in \mathbb{R}_{> 0}$ (set to $1.0$) \\
$E_S$ & Shortcut equilibrium $(u^*(\lambda), 0)$ & $\mathbb{R}_{\ge 0}^2$ \\
$E_C$ & Coherent equilibrium $(u^*(\lambda), v^*(\lambda))$ & $\mathbb{R}_{\ge 0}^2$ \\
\midrule
\multicolumn{3}{l}{\textbf{Coupling Parameters \& Loss Coefficients}} \\
$\lambda$ & Structural compositional pressure & $\lambda \in [0, \infty)$ \\
$\lambda_{\text{crit}}$ & Analytical critical pressure threshold ($b_C / a_C$) & $\lambda_{\text{crit}} \in \mathbb{R}_{> 0}$ (nominal $0.025$) \\
$\hat{\lambda}_{\text{crit}}$ & Empirical fitted critical threshold & $\hat{\lambda}_{\text{crit}} \in [0.015, 0.025]$ \\
$a_S$ & Task alignment gain & $a_S > 0$ (nominal $1.0$) \\
$b_S$ & Shortcut suppression rate per unit pressure & $b_S > 0$ (nominal $1.0$) \\
$a_C$ & Structural coherent gain per unit pressure & $a_C > 0$ (nominal $1.0$) \\
$b_C$ & Intrinsic representational dissipation rate & $b_C > 0$ (nominal $0.025$) \\
$\kappa$ & Non-linear saturation coefficient & $\kappa > 0$ (nominal $1.0$) \\
\midrule
\multicolumn{3}{l}{\textbf{Order Parameters \& Stability Metrics}} \\
$R_0$ & Macroscopic reproductive ratio ($\lambda a_C / b_C$) & $R_0 \in [0, \infty)$ \\
$\mu_{\parallel}$ & Longitudinal Jacobian eigenvalue (along $u$) & $\mu_{\parallel} = -(a_S + \lambda b_S) < 0$ \\
$\mu_{\perp}$ & Transverse Jacobian eigenvalue (along $v$) & $\mu_{\perp} = \lambda a_C - b_C$ \\
$k$ & Separatrix change-point steepness parameter & $k \in \mathbb{R}_{> 0}$ ($k = 79.5$, $95\%$ bootstrap CI $[61.7, 105.0]$, $k_{\text{MLE}} = 93.94$) \\
\multicolumn{3}{l}{\textbf{Representation Geometry \& Curvature Metrics}} \\
$\CKA(X, Y)$ & Linear Centered Kernel Alignment & $\CKA \in [0, 1]$ \\
$g_A^{\text{proj}}(t)$ & Whitened Gradient Cosine Alignment & $g_A^{\text{proj}} \in [-1, 1]$ \\
$P_{\perp}^{\text{emb}}$ & Subspace projection operator for embedding whitening & Orthogonal projector $\in \mathbb{R}^{D \times D}$ \\
$\lambda_{\text{max}}(H_t)$ & Top eigenvalue of loss Hessian matrix $H_t$ & $\lambda_{\text{max}} \in \mathbb{R}_{\ge 0}$ \\
$2/\eta$ & Edge of Stability (EOS) sharpness threshold & $\text{Threshold } = 2000.0$ ($\eta = 0.001$) \\
\bottomrule
\end{tabular}
\end{small}
\end{table}

\section{Mathematical Foundations of the Two-Subspace Flow}
\label{sec:appendix_proof}

In this appendix, we provide the mathematical foundations and derivations for Theorem~\ref{thm:transcritical_bifurcation} concerning the reduced two-subspace continuous gradient flow, including phenomenological parameter identification from projected Hessian traces, Center Manifold perturbation analysis for cross-coupling terms, global Lyapunov stability, and a heuristic preconditioned momentum surrogate.
\subsection{Phenomenological Parameter Identification from Projected Hessian Quadratic Forms and Observable-to-Parameter Bridge}
\label{sec:app_hessian_derivation}
Let $\mathcal{W} \subseteq \mathbb{R}^D$ denote the $D$-dimensional parameter space of the neural network with initialization $w_0 \in \mathcal{W}$. We decompose the parameter space into mutually orthogonal subspaces: $\mathcal{W} = \mathcal{U} \oplus \mathcal{V} \oplus \mathcal{W}_\perp$, where $u \in \mathcal{U} \cong \mathbb{R}_{\ge 0}$ represents the 1D shortcut coordinate, $v \in \mathcal{V} \cong \mathbb{R}_{\ge 0}$ represents the 1D coherent schema coordinate spanning the 2D macroscopic manifold $\mathcal{M} = \mathcal{U} \oplus \mathcal{V}$, and $\mathcal{W}_\perp$ is the $(D-2)$-dimensional orthogonal complement.

Formally, the 1D orthonormal basis vectors spanning $\mathcal{U}$ and $\mathcal{V}$ are defined at initialization $w_0$:
\begin{align}
    \hat{e}_{\mathcal{U}} &\coloneqq \frac{\nabla_w \mathcal{L}_{\text{task}}(w_0)}{\|\nabla_w \mathcal{L}_{\text{task}}(w_0)\|}, \label{eq:app_e_u_def} \\
    \hat{e}_{\mathcal{V}} &\coloneqq \frac{(I_D - P_{\mathcal{U}}) \nabla_w \mathcal{L}_{\text{comp}}(w_0)}{\|(I_D - P_{\mathcal{U}}) \nabla_w \mathcal{L}_{\text{comp}}(w_0)\|}, \label{eq:app_e_v_def}
\end{align}
with orthogonal projection operators $P_{\mathcal{U}} \coloneqq \hat{e}_{\mathcal{U}} \hat{e}_{\mathcal{U}}^\top \in \mathbb{R}^{D \times D}$ and $P_{\mathcal{V}} \coloneqq \hat{e}_{\mathcal{V}} \hat{e}_{\mathcal{V}}^\top \in \mathbb{R}^{D \times D}$, spanning the 2D macroscopic manifold with projector $P_{\mathcal{M}} = P_{\mathcal{U}} + P_{\mathcal{V}}$ ($P_{\mathcal{U}} P_{\mathcal{V}} = 0$). We define $P_{\mathcal{W}_\perp} \coloneqq I_D - P_{\mathcal{M}}$ as the orthogonal projector onto $\mathcal{W}_\perp$, satisfying $P_{\mathcal{U}} + P_{\mathcal{V}} + P_{\mathcal{W}_\perp} = I_D$.

To connect high-dimensional parameter-space trajectories with macroscopic state variables, let $\Delta w(t) = w(t) - w_0$ denote parameter displacement from initialization $w_0 \in \mathcal{W}$. We define the normalized parameter-space projected displacements:
\begin{equation}
    u_{\text{param}}(t) \coloneqq \frac{\| P_{\mathcal{U}} \Delta w(t) \|}{\sqrt{\dim(\mathcal{U})}} = |\hat{e}_{\mathcal{U}}^\top \Delta w(t)|, \quad v_{\text{param}}(t) \coloneqq \frac{\| P_{\mathcal{V}} \Delta w(t) \|}{\sqrt{\dim(\mathcal{V})}} = |\hat{e}_{\mathcal{V}}^\top \Delta w(t)|.
    \label{eq:app_param_displacements}
\end{equation}
In the neighborhood of the nominal shortcut equilibrium $E_S = (\theta_S, 0)$, we define the empirical task loss observable proxy:
\begin{equation}
    u(t) \coloneqq 1 - \frac{\mathcal{L}_{\text{task}}(t)}{\mathcal{L}_{\text{task}}(0)} \in [0, 1].
\end{equation}
Under the simplifying assumptions of first-order gradient alignment along $P_{\mathcal{U}}$, signed projection normalization, and linear functional equivalence, a local second-order Taylor expansion near $E_S$ gives $\mathcal{L}_{\text{task}}(t) \approx \mathcal{L}_{\text{task}}(0)(1 - u_{\text{param}}(t))$, so that $u(t)$ serves as an empirical observable proxy for shortcut reliance. Similarly, for the coherent schema coordinate, linear Centered Kernel Alignment ($\CKA$) between canonical and substituted hidden activations measures representation cosine similarity along the leading structural eigenbasis, defining the proxy:
\begin{equation}
    v(t) \coloneqq \text{CKA}_{\text{linear}}(H_t, H_{\text{perm}}) \cdot \text{AC}(t) \in [0, 1],
\end{equation}
where $\text{AC}(t)$ is the functional Augmentation Consistency. Because CKA is a quadratic statistic invariant under orthogonal transformations, $v(t) \propto v_{\text{param}}(t)$ represents a heuristic local correlation validated by empirical trajectory tracking ($\rho \approx 0.942$), not an analytical equivalence.

Projecting the Hessian matrices evaluated at initialization $w_0$ along the unit basis vectors $\hat{e}_{\mathcal{U}}$ and $\hat{e}_{\mathcal{V}}$, the quadratic forms evaluate to:
\begin{align}
    a_S &\coloneqq \hat{e}_{\mathcal{U}}^\top \nabla_w^2 \mathcal{L}_{\text{task}}(w_0) \hat{e}_{\mathcal{U}} \approx 1.0, \label{eq:app_a_S} \\
    b_C &\coloneqq \hat{e}_{\mathcal{V}}^\top \nabla_w^2 \mathcal{L}_{\text{task}}(w_0) \hat{e}_{\mathcal{V}} \approx 0.025 \pm 0.004, \label{eq:app_b_C} \\
    a_C &\coloneqq \hat{e}_{\mathcal{V}}^\top \left(-\nabla_w^2 \mathcal{L}_{\text{comp}}(w_0)\right) \hat{e}_{\mathcal{V}} \approx 1.0, \label{eq:app_a_C} \\
    b_S &\coloneqq \hat{e}_{\mathcal{U}}^\top \nabla_w^2 \mathcal{L}_{\text{comp}}(w_0) \hat{e}_{\mathcal{U}} \approx 0.5 \text{ to } 1.0. \label{eq:app_b_S}
\end{align}
Here, the unit basis vectors $\hat{e}_{\mathcal{U}}$ and $\hat{e}_{\mathcal{V}}$ are constructed strictly at initialization step 0 prior to training, providing nominal parameter scalings grounded in empirical initialization curvatures. Under standard ERM, the empirical shortcut subspace exhibits steep curvature ($a_S \approx 1.0$) because superficial lexical heuristics align with common training examples. In contrast, the unregularized compositional subspace exhibits very weak quadratic dissipation ($b_C \approx 0.025 \pm 0.004$), reflecting a $40\times$ curvature ratio ($a_S / b_C \approx 40$). Under the compositional penalty $\mathcal{L}_{\text{comp}}$, the structural substitution objective provides balanced unit curvature across both subspaces ($a_C \approx 1.0, b_S \approx 0.5$--$1.0$). Therefore:
\begin{equation}
    \lambda_{\text{crit}} = \frac{b_C}{a_C} = \frac{0.025 \pm 0.004}{1.0} = 0.025 \pm 0.004
\end{equation}
is an analytical property of the loss-landscape Hessian curvature ratio evaluated at step-0 initialization, providing an analytical parameter identification that grounds the nominal continuous bifurcation threshold $\lambda_{\text{crit}} = b_C / a_C \approx 0.025$ in measured task-to-compositional curvature disparity ($a_S / b_C \approx 40$), predicting that the macroscopic phase transition occurs at this identified curvature ratio (empirically measured at $\hat{\lambda}_{\text{crit}} = 0.0238 \in [0.0203, 0.0317]$ on $\hbar$). Crucially, dynamic Lanczos spectral tracking across training trajectories (Section~\ref{sec:hessian_dynamics}) confirms that the projected Hessian curvatures maintain their steep separation ratio $a_S(t) / b_C(t) \gg 1$ throughout optimization until the separatrix is crossed, ensuring that the macroscopic landscape topology remains quantitatively valid throughout training.
\subsection{Heuristic Transverse Dissipation and Dimensional Concentration Analysis}
\label{sec:app_residual_perturbation}
In Section~\ref{sec:subspace_reduction}, empirical dimensionality reduction established that the leading 2-dimensional subspace $\mathcal{U} \oplus \mathcal{V}$ accounts for $86.4\%$ of total representation variance ($D_{\text{eff}} \approx 2.14$). We now analyze parameter perturbations along the $(D-2)$-dimensional orthogonal complement $\mathcal{W}_\perp \coloneqq \mathcal{W} \setminus (\mathcal{U} \oplus \mathcal{V})$ spanning the remaining $13.6\%$ variance.

Let $\mathbf{w}_\perp = P_{\mathcal{W}_\perp} (w - w_0) \in \mathcal{W}_\perp$ denote the transverse deviation vector from the 2-dimensional manifold $\mathcal{M} = \mathcal{U} \oplus \mathcal{V}$. The full continuous gradient flow on $\mathcal{W} \cong \mathbb{R}^D$ is given by:
\begin{equation}
    \dot{w} = -\nabla_w \mathcal{L}(w; \lambda) = -\left( \nabla_w \mathcal{L}_{\text{task}}(w) + \lambda \nabla_w \mathcal{L}_{\text{comp}}(w) + \lambda_{\text{wd}} w \right).
\end{equation}
To address the Hessian sign on the orthogonal complement, we explicitly formulate the subspace decomposition of the compositional penalty Hessian:
\begin{equation}
    \nabla_w^2 \mathcal{L}_{\text{comp}}(w_0) = -a_C P_{\mathcal{V}} + b_S P_{\mathcal{U}} + H_{\mathcal{W}^\perp}^{\text{comp}},
    \label{eq:comp_hessian_decomp}
\end{equation}
where $P_{\mathcal{W}^\perp}^\top H_{\mathcal{W}^\perp}^{\text{comp}} P_{\mathcal{W}^\perp} \succeq 0$ is positive semi-definite on $\mathcal{W}^\perp$. This decomposition separates the negative transverse alignment torque along the 1D coherent coordinate $\mathcal{V}$ (which generates the restorative force $+\lambda a_C v$) from the non-negative structural dissipation on the orthogonal complement $\mathcal{W}^\perp$.

Projecting the dynamics onto the orthogonal subspace $\mathcal{W}_\perp$ via $P_{\mathcal{W}_\perp}$ yields:
\begin{equation}
    \dot{\mathbf{w}}_\perp = -P_{\mathcal{W}_\perp} \nabla_w \mathcal{L}(w; \lambda) = -D_\perp \mathbf{w}_\perp + \mathcal{R}(\mathbf{w}_\perp, u, v),
    \label{eq:w_perp_dynamics}
\end{equation}
where the transverse dissipation operator is defined as:
\begin{equation}
    D_\perp \coloneqq P_{\mathcal{W}_\perp} \left( \nabla_w^2 \mathcal{L}_{\text{task}}(w_0) + \lambda H_{\mathcal{W}^\perp}^{\text{comp}} + \lambda_{\text{wd}} I_D \right) P_{\mathcal{W}_\perp}^\top \in \mathbb{R}^{(D-2) \times (D-2)},
\end{equation}
and $\mathcal{R}(\mathbf{w}_\perp, u, v) = \mathcal{O}(\|\mathbf{w}_\perp\|^2 + \|\mathbf{w}_\perp\| (u + v))$ collects higher-order Taylor residuals.

\begin{lemma}[Strict Dissipativity and Quartic Lyapunov Confinement on Orthogonal Complement]
\label{lem:orthogonal_dissipativity}
For any $\lambda_{\text{wd}} > 0$ and $\lambda \ge 0$, the operator $D_\perp$ is strictly positive definite:
\begin{equation}
    \lambda_{\text{min}}(D_\perp) \ge \lambda_{\text{wd}} > 0.
\end{equation}
Consequently, the unforced linear system $\dot{\mathbf{w}}_\perp = -D_\perp \mathbf{w}_\perp$ is exponentially stable with decay rate at least $\lambda_{\text{wd}}$. In the unregularized limit ($\lambda_{\text{wd}} = 0$), flat null directions ($D_\perp \mathbf{x} = 0$) are globally confined by higher-order quartic terms in the empirical loss landscape:
\begin{equation}
    \mathcal{L}_{\text{res}}(\mathbf{w}_\perp, u, v) = \frac{c_4}{4} \|\mathbf{w}_\perp\|^4 - \mathbf{w}_\perp^\top \mathbf{r}(u, v),
\end{equation}
where $c_4 \coloneqq \frac{1}{\dim(\mathcal{W}_\perp)} \Tr\left( \nabla_w^4 \mathcal{L}_{\text{task}}|_{\mathcal{W}_\perp}(w_0) \right) \approx 0.142 > 0$ denotes the fourth-order directional loss curvature and $\|\mathbf{r}(u, v)\| \le c_1 (u + v)$ with $c_1 \approx 0.05 > 0$. Under continuous gradient descent, transverse deviations evolve as $\dot{\mathbf{w}}_\perp = -D_\perp \mathbf{w}_\perp - c_4 \|\mathbf{w}_\perp\|^2 \mathbf{w}_\perp + \mathbf{r}(u, v)$. Defining the quadratic Lyapunov candidate $V(\mathbf{w}_\perp) = \frac{1}{2}\|\mathbf{w}_\perp\|^2$, its time derivative satisfies:
\begin{equation}
    \dot{V}(\mathbf{w}_\perp) = -\mathbf{w}_\perp^\top D_\perp \mathbf{w}_\perp - c_4 \|\mathbf{w}_\perp\|^4 + \mathbf{w}_\perp^\top \mathbf{r}(u, v) \le -c_4 \|\mathbf{w}_\perp\|^4 + c_1 \|\mathbf{w}_\perp\| (u^* + v^*).
\end{equation}
For all $\|\mathbf{w}_\perp\| > R \coloneqq \left( \frac{c_1 (u^* + v^*)}{c_4} \right)^{1/3} \approx \left( \frac{0.05 \times (1.0 + 0.98)}{0.142} \right)^{1/3} \approx 0.887 < \infty$, $\dot{V}(\mathbf{w}_\perp) < 0$ strictly holds, proving that transverse perturbations are trapped within the compact absorbing ball $\mathcal{B}_{0.887}(0) \coloneqq \{ \mathbf{w}_\perp \in \mathcal{W}_\perp : \|\mathbf{w}_\perp\| \le 0.887 \}$ even in the unregularized limit ($\lambda_{\text{wd}} = 0$).
\end{lemma}

\begin{proof}
Because $\nabla_w^2 \mathcal{L}_{\text{task}}(w_0) \succeq 0$ and $P_{\mathcal{W}_\perp}^\top H_{\mathcal{W}^\perp}^{\text{comp}} P_{\mathcal{W}_\perp} \succeq 0$ are positive semi-definite on $\mathcal{W}_\perp$, for any non-zero $\mathbf{x} \in \mathcal{W}_\perp$:
\begin{equation}
    \mathbf{x}^\top D_\perp \mathbf{x} = \mathbf{x}^\top \nabla_w^2 \mathcal{L}_{\text{task}}(w_0) \mathbf{x} + \lambda \mathbf{x}^\top H_{\mathcal{W}^\perp}^{\text{comp}} \mathbf{x} + \lambda_{\text{wd}} \|\mathbf{x}\|^2 \ge \lambda_{\text{wd}} \|\mathbf{x}\|^2 > 0.
\end{equation}
Thus, all eigenvalues of $-D_\perp$ lie strictly in the left half-plane: $\text{Re}(\nu_i) \le -\lambda_{\text{wd}} < 0$ for all $i \in \{1, \dots, D-2\}$. When $\lambda_{\text{wd}} = 0$, positive semi-definiteness yields $\mathbf{x}^\top D_\perp \mathbf{x} \ge 0$, and the Lyapunov candidate $V(\mathbf{w}_\perp) = \frac{1}{2}\|\mathbf{w}_\perp\|^2$ satisfies $\dot{V}(\mathbf{w}_\perp) \le -c_4 \|\mathbf{w}_\perp\|^4 + c_1 (u^* + v^*) \|\mathbf{w}_\perp\| < 0$ for all $\|\mathbf{w}_\perp\| > R$, establishing ultimate boundedness and radial stability within $\mathcal{B}_R(0)$.
\end{proof}

\begin{conjecture}[Heuristic Transverse Dissipation and Dimensional Concentration Hypothesis]
\label{prop:normal_hyperbolicity}
Under explicit weight decay $\lambda_{\text{wd}} > 0$ and high-order quartic confinement, transverse parameter drift along $\mathcal{W}_\perp$ is bounded within a compact neighborhood of $\mathcal{M} = \mathcal{U} \oplus \mathcal{V}$. In this regime, the leading 2-dimensional continuous flow (Equations~\eqref{eq:ode_u}--\eqref{eq:ode_v}) captures the macroscopic bifurcation structure of the full learning dynamics.
\end{conjecture}

\begin{remark}[Mathematical Status and Open Gaps]
We emphasize that Conjecture~\ref{prop:normal_hyperbolicity} is a phenomenological modelling hypothesis supported by empirical dimensionality concentration ($D_{\text{eff}} \approx 2.14, 86.4\%$ variance) and transverse dissipation scaling. A rigorous proof of full-space invariant manifold persistence would require establishing: (i) the exact existence of a smooth invariant manifold $\mathcal{M}$ on which residual forcing $\mathbf{r}(u, v) \equiv 0$ vanishes identically, (ii) uniform directional coercivity of higher-order terms across all $(D-2)$ orthogonal dimensions (rather than trace-averaged curvature estimates), and (iii) verified normal-hyperbolicity spectral gap conditions across non-compact training domains. We leave these functional-analytic questions as open mathematical problems.
\end{remark}

\subsection{System Equations and Fixed Point Derivation}
Consider the continuous dynamical system defined on the non-negative quadrant $\Omega = \{ (u, v) \in \mathbb{R}^2 : u \ge 0, v \ge 0 \}$:
\begin{align}
    \dot{u} &= f(u, v) = a_S \theta_S - (a_S + \lambda b_S) u, \label{eq:app_u} \\
    \dot{v} &= g(u, v) = (\lambda a_C - b_C) v - \kappa v^2. \label{eq:app_v}
\end{align}
Setting $\dot{u} = 0$ yields a unique $u$-coordinate for any fixed point:
\begin{equation}
    u^*(\lambda) = \frac{a_S \theta_S}{a_S + \lambda b_S}.
\end{equation}
Since $a_S, \theta_S, b_S > 0$ and $\lambda \ge 0$, $u^*(\lambda)$ is strictly positive, finite, and monotonically decreasing with $\lambda$.

Setting $\dot{v} = 0$ yields:
Setting $\dot{v} = 0$ yields $v_1^* = 0$ or $v_2^* = \frac{\lambda a_C - b_C}{\kappa}$. Formulated on the physical non-negative state space $\Omega = \mathbb{R}_{\ge 0}^2$, the bifurcation is a boundary equilibrium bifurcation:
\begin{enumerate}[leftmargin=*,noitemsep]
    \item The shortcut equilibrium $E_S = (u^*(\lambda), 0) \in \{v=0\}$ exists on the invariant boundary for all $\lambda \ge 0$.
    \item For subcritical pressure $\lambda < \lambda_{\text{crit}} \coloneqq b_C / a_C$, $v_2^*(\lambda) < 0$, so the non-trivial equilibrium $E_C = \left( u^*(\lambda), \frac{\lambda a_C - b_C}{\kappa} \right)$ resides outside $\Omega$ in the unphysical lower half-plane ($\mathbb{R} \times \mathbb{R}_{<0}$), leaving the boundary fixed point $E_S$ as the unique globally asymptotically stable sink on the closed quadrant $\Omega = \mathbb{R}_{\ge 0}^2$.
    \item At $\lambda = \lambda_{\text{crit}}$, $E_C$ collides with $E_S$ at $(u^*(\lambda_{\text{crit}}), 0)$ on the invariant boundary $\{v=0\}$.
    \item For supercritical pressure $\lambda > \lambda_{\text{crit}}$, $E_C$ crosses into the physical interior $\Omega^\circ = \mathbb{R}_{> 0}^2$ ($v_2^*(\lambda) > 0$), exchanging stability with $E_S$ such that $E_S$ becomes an unstable saddle and $E_C$ emerges as the unique globally asymptotically stable sink on the physical interior $\Omega^\circ$ (and for all initial states with $v(0) > 0$).
\end{enumerate}

\subsection{Linearization and Jacobian Spectrum}
The Jacobian matrix of the vector field $(f, g)^T$ at an arbitrary state $(u, v)$ is given by:
\begin{equation}
    J(u, v) = \begin{pmatrix}
        \frac{\partial f}{\partial u} & \frac{\partial f}{\partial v} \\
        \frac{\partial g}{\partial u} & \frac{\partial g}{\partial v}
    \end{pmatrix} = \begin{pmatrix}
        -(a_S + \lambda b_S) & 0 \\
        0 & (\lambda a_C - b_C) - 2\kappa v
    \end{pmatrix}.
\end{equation}
Because $J(u, v)$ is strictly diagonal for all $(u, v)$, the eigenvectors are aligned with the coordinate axes: $\mathbf{e}_u = (1, 0)^T$ (longitudinal shortcut direction) and $\mathbf{e}_v = (0, 1)^T$ (transverse coherent direction).

\paragraph{Case 1: Subcritical Regime ($\lambda < \lambda_{\text{crit}}$).}
Evaluating $J$ at $E_S$:
\begin{equation}
    J(E_S) = \begin{pmatrix}
        -(a_S + \lambda b_S) & 0 \\
        0 & \lambda a_C - b_C
    \end{pmatrix}.
\end{equation}
The eigenvalues are:
\begin{equation}
    \mu_{\parallel} = -(a_S + \lambda b_S) < 0, \quad \mu_{\perp} = \lambda a_C - b_C < 0.
\end{equation}
Since both eigenvalues are strictly negative and real, $E_S$ is a stable hyperbolic sink.

\paragraph{Case 2: Bifurcation Point ($\lambda = \lambda_{\text{crit}} = b_C / a_C$).}
At $\lambda = \lambda_{\text{crit}}$, $\mu_{\parallel} = -(a_S + \lambda_{\text{crit}} b_S) < 0$ and $\mu_{\perp} = 0$. The system possesses a non-hyperbolic fixed point at $E_S$, with a 1-dimensional center manifold tangent to $\mathbf{e}_v$.

\paragraph{Case 3: Supercritical Regime ($\lambda > \lambda_{\text{crit}}$).}
Evaluating $J$ at $E_S$:
\begin{equation}
    \mu_{\parallel} = -(a_S + \lambda b_S) < 0, \quad \mu_{\perp} = \lambda a_C - b_C > 0.
\end{equation}
Thus, $E_S$ becomes a hyperbolic saddle point, unstable along the transverse coherent direction $\mathbf{e}_v$.

Evaluating $J$ at $E_C = \left( u^*(\lambda), \frac{\lambda a_C - b_C}{\kappa} \right)$:
\begin{equation}
    J(E_C) = \begin{pmatrix}
        -(a_S + \lambda b_S) & 0 \\
        0 & (\lambda a_C - b_C) - 2\kappa \left( \frac{\lambda a_C - b_C}{\kappa} \right)
    \end{pmatrix} = \begin{pmatrix}
        -(a_S + \lambda b_S) & 0 \\
        0 & -(\lambda a_C - b_C)
    \end{pmatrix}.
\end{equation}
The eigenvalues are:
\begin{equation}
    \mu_{\parallel} = -(a_S + \lambda b_S) < 0, \quad \mu_{\perp} = -(\lambda a_C - b_C) < 0.
\end{equation}
Both eigenvalues are strictly negative, proving that $E_C$ is a locally asymptotically stable sink.

\subsection{Center Manifold and Perturbation Analysis for Cross-Coupling Terms}
\label{sec:app_center_manifold}
We now extend the theoretical framework to non-separable loss landscapes incorporating cross-subspace couplings. We analyze two complementary formulations: (i) physical curvature-modulating coupling where shortcut reliance modulates transverse schema curvature, and (ii) general bilinear coupling around the perturbed equilibrium.

\paragraph{1. Physical Curvature-Modulating Coupling $\frac{1}{2}\gamma u v^2$.}
In multi-head self-attention and MLP feed-forward layers, early lexical shortcut heads activate rapidly ($u \to \theta_S$), which projects high-variance representations that increase the effective loss curvature of downstream compositional query-key projections ($b_C + \gamma u$). This flattens the transverse compositional gradient until supercritical pressure $\lambda > \lambda_{\text{crit}}$ counterbalances the shortcut-induced dissipation. We formalize this via the physical coupled task loss:
\begin{equation}
    \mathcal{L}_{\text{task}}(u, v) = \frac{1}{2} a_S (\theta_S - u)^2 + \frac{1}{2} (b_C + \gamma u) v^2 + \frac{1}{3} \kappa v^3,
\end{equation}
where $\gamma > 0$ denotes the cross-subspace curvature modulation coefficient. Under continuous gradient flow with the compositional penalty $\mathcal{L}_{\text{comp}}(u, v) = \frac{1}{2} b_S u^2 - \frac{1}{2} a_C v^2$, the dynamical system becomes:
\begin{align}
    \dot{u} &= a_S \theta_S - (a_S + \lambda b_S) u - \frac{1}{2} \gamma v^2, \label{eq:curv_coupled_u} \\
    \dot{v} &= \left(\lambda a_C - (b_C + \gamma u)\right) v - \kappa v^2. \label{eq:curv_coupled_v}
\end{align}
Crucially, evaluating along the shortcut axis $v = 0$ yields $\dot{v} = 0$ identically for all $u$, proving that the invariant boundary manifold $\{v = 0\}$ is strictly preserved for all $\lambda \ge 0$. Consequently, the shortcut equilibrium:
\begin{equation}
    E_S = (u^*(\lambda), 0) = \left( \frac{a_S \theta_S}{a_S + \lambda b_S}, 0 \right)
\end{equation}
remains an exact fixed point for all $\lambda \ge 0$. Linearizing around $E_S$, the transverse eigenvalue is $\mu_\perp(\lambda) = \lambda a_C - \left(b_C + \gamma u^*(\lambda)\right)$. The critical bifurcation threshold $\tilde{\lambda}_{\text{crit}}$ is the unique solution to $\mu_\perp(\tilde{\lambda}_{\text{crit}}) = 0$:
\begin{equation}
    \tilde{\lambda}_{\text{crit}} = \frac{b_C + \gamma u^*(\tilde{\lambda}_{\text{crit}})}{a_C} = \frac{b_C + \gamma \frac{a_S \theta_S}{a_S + \tilde{\lambda}_{\text{crit}} b_S}}{a_C}.
\end{equation}
To determine the reduced flow on the center manifold at $\lambda \approx \tilde{\lambda}_{\text{crit}}$, let $\tilde{u} \coloneqq u - u^*(\lambda)$ and $A \coloneqq a_S + \lambda b_S > 0$. The deviation dynamics satisfy $\dot{\tilde{u}} = -A \tilde{u} - \frac{1}{2} \gamma v^2$. Because the linear eigenvalue $-A < 0$ is strictly negative, the center manifold $\tilde{u} = h(v)$ is given by:
\begin{equation}
    h(v) = -\frac{\gamma}{2A} v^2 + \mathcal{O}(v^3).
\end{equation}
Substituting $\tilde{u} = h(v)$ into Equation~\eqref{eq:curv_coupled_v} yields the reduced flow on the 1D center manifold:
\begin{align}
    \dot{v} &= \left( \lambda a_C - b_C - \gamma (u^*(\lambda) + h(v)) \right) v - \kappa v^2 \nonumber \\
    &= \tilde{\mu}_\perp(\lambda) v - \kappa v^2 - \gamma h(v) v \nonumber \\
    &= \tilde{\mu}_\perp(\lambda) v - \kappa v^2 + \frac{\gamma^2}{2A} v^3 + \mathcal{O}(v^4),
\end{align}
where $\tilde{\mu}_\perp(\lambda) \coloneqq \lambda a_C - (b_C + \gamma u^*(\lambda))$. Notice that the leading non-linear term remains $-\kappa v^2$ (with $\kappa > 0$), while the cross-coupling introduces a cubic perturbation $+\frac{\gamma^2}{2A} v^3$. Because the quadratic term dominates near $v = 0$, the transcritical normal form is structurally preserved without altering the bifurcation type.

\paragraph{2. General Bilinear Coupling $\gamma u v$ and Imperfect Bifurcation.}
For a general bilinear coupling term $\mathcal{L}_{\text{task}}(u, v) = \frac{1}{2} a_S (\theta_S - u)^2 + \frac{1}{2} b_C v^2 + \frac{1}{3} \kappa v^3 + \gamma u v$ with $|\gamma| < \gamma_{\text{max}} \coloneqq \sqrt{a_S b_C} = \sqrt{1.0 \times 0.025} \approx 0.1581$, the gradient flow yields:
\begin{align}
    \dot{u} &= a_S \theta_S - (a_S + \lambda b_S) u - \gamma v, \label{eq:gen_coupled_u} \\
    \dot{v} &= (\lambda a_C - b_C) v - \kappa v^2 - \gamma u. \label{eq:gen_coupled_v}
\end{align}
Notice that the non-vanishing $-\gamma u$ term in Equation~\eqref{eq:gen_coupled_v} breaks the strict invariance of the boundary manifold $\{v=0\}$, as $\dot{v}|_{v=0} = -\gamma u \ne 0$. For $\gamma > 0$ and $u > 0$, the vector field at $v = 0$ points outside the non-negative quadrant $\mathbb{R}_{\ge 0}^2$. Consequently, generic bilinear perturbations transform the transcritical bifurcation into an imperfect bifurcation with an avoided crossing of boundary layer width $\mathcal{O}(\gamma)$. Setting $(\dot{u}, \dot{v}) = (0, 0)$ defines the perturbed fixed point $(u^*_\gamma, v^*_\gamma) = (u^*(\lambda) + \mathcal{O}(\gamma), \mathcal{O}(\gamma))$. Expanding the vector field in deviations $\delta u = u - u^*_\gamma$ and $\delta v = v - v^*_\gamma$, the Jacobian matrix is:
\begin{equation}
    J_\gamma = \begin{pmatrix}
        -(a_S + \lambda b_S) & -\gamma \\
        -\gamma & (\lambda a_C - b_C) - 2\kappa v^*_\gamma
    \end{pmatrix}.
\end{equation}
The trace and determinant evaluate to $\Tr(J_\gamma) = -(a_S + \lambda b_S) + (\lambda a_C - b_C - 2\kappa v^*_\gamma)$ and $\det(J_\gamma) = -(a_S + \lambda b_S)(\lambda a_C - b_C - 2\kappa v^*_\gamma) - \gamma^2$. For all $|\gamma| < \gamma_{\text{max}} \coloneqq \sqrt{a_S b_C} \approx 0.1581$, the determinant crosses zero at $\tilde{\lambda}_{\text{crit}} = \frac{b_C + 2\kappa v^*_\gamma - \gamma^2 / (a_S + \tilde{\lambda}_{\text{crit}} b_S)}{a_C} \approx \lambda_{\text{crit}}$, where the transverse eigenvalue undergoes a real stability exchange from negative to positive. This confirms that generic bilinear perturbations smooth the transcritical crossing into an avoided crossing, while preserving the macroscopic stability exchange within an $\mathcal{O}(\gamma)$ neighborhood. Strict transcritical bifurcations on $\mathbb{R}_{\ge 0}^2$ require symmetry-preserving perturbations respecting $\{v=0\}$ boundary invariance.
\subsection{Global Asymptotic Stability via Lyapunov Functions and Critical Solution}
We now analyze global stability on $\Omega$. Notice that Equation~\eqref{eq:app_u} is completely decoupled from $v$ and linear:
\begin{equation}
    u(t) = u^*(\lambda) + \left( u(0) - u^*(\lambda) \right) e^{-(a_S + \lambda b_S) t} \implies \lim_{t \to \infty} u(t) = u^*(\lambda),
\end{equation}
with exponential convergence rate $a_S + \lambda b_S > 0$.

For the coherent coordinate $v(t)$, Equation~\eqref{eq:app_v} is a classical Bernoulli ODE:
\begin{equation}
    \dot{v} = \alpha v - \kappa v^2, \quad \text{where } \alpha \coloneqq \lambda a_C - b_C.
\end{equation}
We consider three cases based on the sign of $\alpha$:
\begin{enumerate}[leftmargin=*]
    \item \textbf{Critical Point ($\alpha = 0 \iff \lambda = \lambda_{\text{crit}}$):}
    At the critical bifurcation boundary, the linear term vanishes identically, giving:
    \begin{equation}
        \dot{v} = -\kappa v^2 \implies \frac{dv}{v^2} = -\kappa \, dt \implies -\frac{1}{v(t)} + \frac{1}{v(0)} = -\kappa t.
    \end{equation}
    Solving for $v(t)$ yields the exact algebraic decay solution:
    \begin{equation}
        v(t) = \frac{v(0)}{1 + \kappa v(0) t} = \mathcal{O}(t^{-1}) \quad \text{as } t \to \infty.
        \label{eq:critical_ode_solution}
    \end{equation}
    This proves that at the critical boundary, exponential convergence is replaced by algebraic decay (critical slowing down).
    \item \textbf{Subcritical ($\alpha < 0 \iff \lambda < \lambda_{\text{crit}}$):}
    Applying the substitution $z(t) = 1/v(t)$ for $v(0) > 0$ yields $\dot{z} = -\alpha z + \kappa$, with solution:
    \begin{equation}
        v(t) = \frac{1}{\frac{\kappa}{\alpha} + \left( \frac{1}{v(0)} - \frac{\kappa}{\alpha} \right) e^{-\alpha t}}.
    \end{equation}
    Since $\alpha < 0$, $e^{-\alpha t} = e^{|\alpha| t} \to \infty$ as $t \to \infty$, which implies $\lim_{t \to \infty} v(t) = 0$. Hence $(u(t), v(t)) \to E_S$ for all $(u(0), v(0)) \in \Omega = \mathbb{R}_{\ge 0}^2$.
    \item \textbf{Supercritical ($\alpha > 0 \iff \lambda > \lambda_{\text{crit}}$):}
    Since $\alpha > 0$, $e^{-\alpha t} \to 0$ as $t \to \infty$, which implies $\lim_{t \to \infty} v(t) = \frac{\alpha}{\kappa} = \frac{\lambda a_C - b_C}{\kappa} = v^*$. Hence $(u(t), v(t)) \to E_C$ for all initial states in the physical interior $\Omega^\circ = \mathbb{R}_{> 0}^2$ (all initial states with $v(0) > 0$). Trajectories on the invariant boundary $v(0) = 0$ remain at $v(t) = 0$ and converge to $E_S$.
\end{enumerate}

Furthermore, on the physical interior $\Omega^\circ = \mathbb{R}_{> 0}^2$, we define the strict Lyapunov candidate function:
\begin{equation}
    V(u, v) = \frac{1}{2} (u - u^*(\lambda))^2 + \int_{v^*(\lambda)}^v \frac{s - v^*(\lambda)}{s} \, ds.
\end{equation}
Notice that $V(u, v)$ possesses a logarithmic singularity at $v = 0$ ($\int_{v^*}^v \frac{s - v^*}{s} ds = v - v^* - v^* \ln(v/v^*)$) and is therefore defined strictly on $\mathbb{R}_{\ge 0} \times \mathbb{R}_{> 0}$. Its time derivative along trajectories evaluates to:
\begin{equation}
    \dot{V}(u, v) = -(a_S + \lambda b_S)(u - u^*)^2 - \kappa (v - v^*)^2 \le 0,
\end{equation}
with $\dot{V} = 0$ if and only if $(u, v) = E_C$, establishing global asymptotic stability on $\Omega^\circ$ via LaSalle's Invariance Principle.
\subsection{Phenomenological Stochastic Escape Hypothesis and Finite-Horizon Baseline Analysis}
\label{sec:app_langevin_escape}
To conceptually explore the discrete mini-batch optimization observations, we formulate discrete stochastic gradient descent as a continuous Langevin stochastic differential equation (SDE):
\begin{equation}
    d w_t = -\nabla_w \mathcal{L}(w_t; \lambda) \, dt + \sqrt{2 D} \, dW_t,
    \label{eq:langevin_sde}
\end{equation}
where $W_t$ is a standard $D$-dimensional Brownian motion, and $D \coloneqq \frac{\eta}{2 B} \Sigma_{\text{batch}}$ represents the effective anisotropic diffusion tensor induced by the mini-batch gradient covariance:
\begin{equation}
    \Sigma_{\text{batch}} \coloneqq \mathbb{E}_{B \sim \mathcal{D}} \left[ \left( \nabla \mathcal{L}_B(w) - \nabla \mathcal{L}(w) \right) \left( \nabla \mathcal{L}_B(w) - \nabla \mathcal{L}(w) \right)^\top \right],
\end{equation}
under learning rate $\eta$ and mini-batch size $B$.

In the deterministic 2D reduced ODE (Theorem~\ref{thm:transcritical_bifurcation}), the subcritical regime ($\lambda < \lambda_{\text{crit}} = b_C / a_C$) features $E_S = (u^*(\lambda), 0)$ as a globally asymptotically stable sink on $\Omega$, with $\dot{v} < 0$ for all $v > 0$. Crucially, the deterministic 2D reduced ODE possesses no interior saddle point or potential barrier. The minor baseline subcritical escape ($2/30 = 6.7\%$ of seeds on $\hbar$) observed under discrete mini-batch optimization is an empirical baseline finding arising from finite-horizon stochastic fluctuations over unmodeled high-dimensional features.

We frame this conceptually via a \textbf{qualitative phenomenological modeling hypothesis} (Level~3 open-scope): if effective stochastic fluctuations on the full loss landscape interact with unmodeled transverse degrees of freedom, an effective barrier scale $\Delta V$ and transverse diffusion coefficient $D_v \coloneqq \mathbf{e}_v^\top D \mathbf{e}_v = \frac{\eta}{2 B} \sigma_v^2$ would yield a qualitative Kramers escape rate \citep{Kramers1940Brownian,Gardiner2009Stochastic}:
\begin{equation}
    \Gamma_{\text{escape}} \propto \exp\left( -\frac{\Delta V}{D_v} \right).
    \label{eq:kramers_rate}
\end{equation}
In the continuous deterministic limit ($B \to \infty$ or $\eta \to 0$), the diffusion tensor vanishes ($D_v \to 0$), so that $\lim_{D_v \to 0} \Gamma_{\text{escape}} = 0$, recovering the deterministic result that $E_S$ is strictly absorbing (Theorem~\ref{thm:transcritical_bifurcation}).

\paragraph{Initial Exploration and 20k-Step Anti-Grokking Stability.}
A key empirical observation is that the minor baseline escape fraction ($2/30 = 6.7\%$ over $T=2000$ steps) does not lead to delayed spontaneous escape (``grokking'') when training is extended to $T = 20{,}000$ steps (Section~\ref{sec:extended_horizon_grokking}). Mechanistically, as training proceeds under standard ERM, shortcut feature specialization deepens ($u(t) \to \theta_S$), which drives empirical task loss to near-zero levels. This causes mini-batch gradient variance to diminish ($\Sigma_{\text{batch}}(t) \to 0$), reducing the effective diffusion coefficient $D_v(t) \to 0$. As the noise-to-barrier ratio $D_v(t) / \Delta V \to 0$, stochastic escape is exponentially suppressed ($\Gamma_{\text{escape}}(t) \to 0$). Consequently, subcritical models that fail to escape within early exploratory iterations ($t \le 500$) become permanently trapped in the shortcut basin $E_S$, explaining the absence of delayed grokking across $20{,}000$ steps.
\subsection{Heuristic Heavy-Ball Surrogate with Scalar Preconditioning}
\label{sec:app_adamw_bridge}
Practical neural networks are trained with discrete-time optimizers incorporating momentum and adaptive learning rates, such as AdamW \citep{Loshchilov2019Decoupled}. Let $w_t \in \mathbb{R}^D$ denote the parameters at discrete step $t$. The AdamW update equations with first-moment coefficient $\beta_1 \in [0.9, 0.99]$, second-moment coefficient $\beta_2 \in [0.98, 0.999]$, and learning rate $\eta > 0$ are:
\begin{align}
    m_t &= \beta_1 m_{t-1} + (1 - \beta_1) \nabla_w \mathcal{L}(w_t), \\
    v_t &= \beta_2 v_{t-1} + (1 - \beta_2) (\nabla_w \mathcal{L}(w_t))^2, \\
    w_{t+1} &= w_t - \eta \frac{m_t}{\sqrt{v_t} + \epsilon} - \eta \lambda_{\text{wd}} w_t.
\end{align}
To analyze how momentum and preconditioning qualitatively interact with the vector field, we study an idealized continuous-time surrogate model. In the limit $\eta \to 0$ with effective momentum timescale $\tau_m = \frac{\eta}{1 - \beta_1}$, the dynamics follow the idealized coupled second-order system:
\begin{align}
    \ddot{w} + \frac{1 - \beta_1}{\eta} \dot{w} &= -\hat{D}_t^{-1/2} \nabla_w \mathcal{L}(w), \label{eq:underdamped_adamw} \\
    \dot{v}_t &= \frac{1 - \beta_2}{\eta} \left( (\nabla_w \mathcal{L}(w))^2 - v_t \right),
\end{align}
where $\hat{D}_t \coloneqq \diag(\sqrt{v_t} + \epsilon) \in \mathbb{R}^{D \times D}$ is the diagonal preconditioner.

\paragraph{Preconditioned Heavy-Ball Model Analysis.}
Because $\epsilon > 0$ and gradient norms are bounded on compact trajectories, the preconditioner satisfies uniform spectral bounds:
\begin{equation}
    \epsilon I_D \preceq \hat{D}_t \preceq M I_D, \quad \text{for some } M < \infty.
\end{equation}
Within this idealized surrogate, $\hat{D}_t^{-1/2}$ acts as a smooth, positive-definite Riemannian metric tensor $G^{-1}(w)$ that rescales velocity along each coordinate without altering the directional sign of the gradient flow.
Linearizing Equation~\eqref{eq:underdamped_adamw} along the transverse coherent coordinate $v$ near the shortcut equilibrium $E_S$ gives:
\begin{equation}
    \ddot{v} + \Gamma \dot{v} - \hat{D}_v^{-1/2} \mu_{\perp} v = 0, \quad \text{where } \Gamma \coloneqq \frac{1 - \beta_1}{\eta} > 0, \quad \mu_{\perp} = \lambda a_C - b_C,
\end{equation}
and $\hat{D}_v^{-1/2} > 0$ is the scalar preconditioning factor along $\mathcal{V}$.
The characteristic equation is $r^2 + \Gamma r - \hat{D}_v^{-1/2} \mu_{\perp} = 0$, with roots:
\begin{equation}
    r_{\pm} = -\frac{\Gamma}{2} \pm \sqrt{\frac{\Gamma^2}{4} + \hat{D}_v^{-1/2} \mu_{\perp}}.
\end{equation}
\begin{itemize}[leftmargin=*]
    \item For $\lambda < \lambda_{\text{crit}}$, $\mu_{\perp} < 0$, yielding $\frac{\Gamma^2}{4} + \hat{D}_v^{-1/2} \mu_{\perp} < \frac{\Gamma^2}{4}$. Because $\hat{D}_v^{-1/2} > 0$, both roots have strictly negative real parts ($\text{Re}(r_{\pm}) < 0$), and $v(t) \to 0$ exponentially.
    \item For $\lambda > \lambda_{\text{crit}}$, $\mu_{\perp} > 0$, yielding a unique strictly positive real eigenvalue $r_+ > 0$. In the high-friction limit ($\Gamma \gg 1$), Taylor expansion yields:
    \begin{equation}
        r_+ = -\frac{\Gamma}{2} + \frac{\Gamma}{2} \sqrt{1 + \frac{4 \hat{D}_v^{-1/2} \mu_{\perp}}{\Gamma^2}} \approx \frac{\hat{D}_v^{-1/2} \mu_{\perp}}{\Gamma} > 0.
    \end{equation}
\end{itemize}
In this simplified surrogate, the sign of the eigenvalue satisfies $\text{sign}(\hat{D}_v^{-1/2} \mu_\perp) = \text{sign}(\mu_\perp)$, indicating that the zero-crossing of the transverse stability eigenvalue occurs at $\mu_\perp = 0 \iff \lambda = \lambda_{\text{crit}} = b_C / a_C$.

\paragraph{Omitted AdamW Mechanisms and Scope Limitations.}
We explicitly emphasize that this continuous heavy-ball surrogate does NOT constitute a mathematical proof of discrete AdamW convergence or exact threshold preservation. Real discrete AdamW differs from the surrogate through several coupled mechanisms:
\begin{enumerate}[leftmargin=*,noitemsep]
    \item Discrete-time step updates with finite step size $\eta > 0$;
    \item Coupled non-stationary first- and second-moment state variables ($m_t, v_t$);
    \item Coordinate-dependent non-uniform scaling across the $D$-dimensional parameter space;
    \item Bias correction factors ($1 - \beta_1^t, 1 - \beta_2^t$); and
    \item Decoupled weight-decay interactions ($-\eta \lambda_{\text{wd}} w_t$).
\end{enumerate}
The surrogate serves strictly as an illustrative conceptual model showing that scalar metric preconditioning does not invert the sign of the transverse restoring force.

\paragraph{Transient vs.\ Steady-State Momentum Validity.}
For standard training hyperparameters ($\beta_1 = 0.9, \eta = 10^{-3}$), the effective friction coefficient is $\Gamma = (1 - 0.9)/10^{-3} = 100.0 \gg 1$. The inertial relaxation timescale is $\tau_m \coloneqq \frac{\eta}{1 - \beta_1} = 0.01$ ($\equiv 10$ discrete gradient steps), which is more than an order of magnitude faster than the empirical takeoff latency ($\hat{\tau} \approx 148$ steps, Section~\ref{sec:rate_ordering}). Consequently, during early training transients ($t < \hat{\tau}$), momentum acts as a velocity accelerator along the unstable manifold without altering the nominal sign of $\mu_{\perp} = \lambda a_C - b_C$.

\paragraph{Critical Discretization Breakdown Limit.}
Under discrete Euler/gradient updates $w_{t+1} = w_t - \eta \tilde{\nabla} \mathcal{L}(w_t)$, linear stability around local attractors requires $|1 - \eta \lambda_i(H_t)| < 1$, deriving the classical step-size stability boundary:
\begin{equation}
    \eta_{\text{crit}} \coloneqq \frac{2}{\lambda_{\text{max}}(H_t)}.
\end{equation}
Across all 960 production runs and sampled Hessian checkpoints (yielding $\lambda_{\text{max}} \approx (2.23 \pm 3.97) \times 10^{-4}$ for Transformer 2L with peak $0.00205$, and transient exploratory peak $\approx 2.071$ observed during initial exploratory screening before settling to sub-unity curvature; Section~\ref{sec:hessian_dynamics}), the sharpness ceiling is $2/\eta = 2000.0 \gg \lambda_{\text{max}}$, yielding critical learning rate $\eta_{\text{crit}} = \frac{2}{\lambda_{\text{max}}} \gg \eta = 1.0 \times 10^{-3}$. Empirical Lanczos measurements confirm that training operates far below the plain-gradient-descent Euler instability threshold ($2/\eta = 2000.0$), ruling out chaotic discretization explosion. This empirical EOS bound does not prove exact ODE trajectory tracking for full discrete AdamW optimization.
\section{ODE Solver Implementation \& Numerical Scheme}
\label{sec:appendix_solver}

To verify the analytical predictions numerically and generate continuous reference trajectories, we implement an adaptive Runge-Kutta Dormand-Prince (RK45) integrator with stiffness detection in \texttt{paper/src/continuous/solver.py} and \texttt{two\_subspace\_ode.py}.

\subsection{Numerical Integration Scheme}
The system of ODEs in Equations~\eqref{eq:ode_u}--\eqref{eq:ode_v} is integrated on $t \in [0, 50.0]$ with adaptive step size:
\begin{itemize}[noitemsep]
    \item Absolute error tolerance: $\text{atol} = 10^{-8}$.
    \item Relative error tolerance: $\text{rtol} = 10^{-8}$.
    \item Maximum step size: $\Delta t_{\text{max}} = 0.05$.
\end{itemize}

Stiffness is monitored via the local spectral radius $\rho(J) = \max\{ |a_S + \lambda b_S|, |(\lambda a_C - b_C) - 2\kappa v| \}$. Because the maximum local spectral radius satisfies $\rho(J) \le 2.0$ across all explored parameter ranges, the maximum step size $\Delta t_{\text{max}} = 0.05$ satisfies $\rho(J) \Delta t_{\text{max}} \le 0.1$, placing the integration safely within the linear stability region of explicit RK45.

\section{Benchmark Grammar Formalization \& Dataset Cards}
\label{sec:appendix_datasets}

\subsection{Homomorphic Command Algebra (\texorpdfstring{$\hbar$}{H-Bar} Benchmark Specification)}
The $\hbar$ (H-Bar) benchmark is a closed, synthetic formal algebra designed to evaluate compositional systematicity in a mathematically rigorous environment with provable support disjointness.

\paragraph{Formal Grammar Definition.}
Let $\mathcal{G}_{\hbar} = (\mathcal{V}_{\text{nt}}, \mathcal{V}_{\text{t}}, \mathcal{R}, S)$ denote the context-free command grammar, where:
\begin{itemize}[noitemsep]
    \item Non-terminals: $\mathcal{V}_{\text{nt}} = \{S, E, P, \text{Op}\}$.
    \item Terminals: $\mathcal{V}_{\text{t}} = \{a, b, c, d\} \cup \{\circ, \star, \text{inv}, \text{rev}, (, )\}$.
    \item Production rules $\mathcal{R}$:
    \begin{align}
        S &\to E \\
        E &\to E \circ E \mid E \star E \mid \text{inv}(E) \mid \text{rev}(E) \mid P \\
        P &\to a \mid b \mid c \mid d
    \end{align}
\end{itemize}

\paragraph{Homomorphic Target Execution Semantics.}
The evaluation map $\phi: \mathcal{L}(\mathcal{G}_{\hbar}) \to \mathcal{A}$ maps command strings to target execution traces via a strict homomorphism:
\begin{equation}
    \phi(E_1 \circ E_2) = \phi(E_1) \cdot \phi(E_2), \quad \phi(E_1 \star E_2) = \text{interleave}(\phi(E_1), \phi(E_2)), \quad \phi(\text{inv}(E)) = \overline{\phi(E)}.
\end{equation}

\paragraph{Dataset Partitions and Support Disjointness.}
The dataset is partitioned into three splits:
\begin{itemize}[noitemsep]
    \item \textbf{Training Set ($N_{\text{train}} = 10{,}000$):} Sequences generated with tree depth $d \in \{1, 2, 3\}$, source length $L_{\text{src}} \in [4, 16]$, target length $L_{\text{tgt}} \in [2, 8]$. Primitives $a, b, c, d$ appear in limited, co-occurring pairs.
    \item \textbf{Validation Set ($N_{\text{val}} = 1{,}000$):} In-distribution held-out split with depth $d \le 3$.
    \item \textbf{Out-of-Distribution Test Set ($N_{\text{test}} = 1{,}500$):} Deep recursive compositional split with tree depth $d \in \{4, 5\}$, source length $L_{\text{src}} \in [12, 24]$, target length $L_{\text{tgt}} \in [6, 12]$.
\end{itemize}

\paragraph{Cryptographic Support Disjointness Proof.}
Let $\text{Trees}(d)$ denote the set of abstract syntax trees of depth $d$. By construction:
\begin{equation}
    \text{supp}(\mathcal{D}_{\text{train}}) \subseteq \bigcup_{k=1}^3 \text{Trees}(k), \quad \text{supp}(\mathcal{D}_{\text{test}}) \subseteq \bigcup_{k=4}^5 \text{Trees}(k) \implies \text{supp}(\mathcal{D}_{\text{train}}) \cap \text{supp}(\mathcal{D}_{\text{test}}) = \emptyset.
\end{equation}
Because zero subtrees of depth $\ge 4$ appear in the training split, zero-shot OOD accuracy measures pure systematic generalization rather than memorized sub-expression recall.

\subsection{Natural Language Compositional Benchmarks \& Cryptographic Test Isolation}
\label{sec:appendix_external_benchmarks}
To evaluate whether the phase-boundary phenomenon generalizes to established linguistic benchmarks, we ingest three widely studied community suites with strict syntactic and derivational test isolation:
\begin{itemize}[leftmargin=*,noitemsep]
    \item \textbf{SCAN (\texttt{add\_primitive\_jump}):} 14,670 training examples containing primitive \texttt{jump} only in isolated atomic contexts (\texttt{"jump"} $\to$ \texttt{I\_JUMP}); 7,706 test examples requiring composition of \texttt{jump} with complex modifiers (\texttt{"jump around left twice"} $\to$ sequences of 16 actions) \citep{Lake2018Generalization}. Substitution pairs are derived strictly from the training primitive lexicon.
    \item \textbf{COGS / ReCOGS:} 24,155 training examples and 21,000 structural test examples evaluating novel syntactic role assignment (agent vs.\ theme transpositions, passive nominalizations, and recursive prepositional phrase attachments) \citep{Kim2020COGS,Wu2023ReCOGS}. Substitution pairs are constructed from training grammar production rules.
    \item \textbf{PCFG-SET:} 95,000 context-free grammar sentences evaluated on systematicity and substitutivity splits with sequence lengths up to 30 tokens \citep{Hupkes2019Compositionality}. Substitution pairs are drawn from training non-terminal expansions.
\end{itemize}

\paragraph{Deterministic Test-Isolation Guarantee.}
All four benchmark suites were audited for data contamination and split overlap using the automated leakage audit suite (\texttt{paper/src/data/audit\_leakage.py}). Pairwise evaluation confirms strict support disjointness across training, validation, and test partitions ($\text{supp}(\mathcal{D}_{\text{train}}) \cap \text{supp}(\mathcal{D}_{\text{test}}) = \emptyset$, with exactly 0 exact matches, 0 derivational leaks, and $0.0\%$ leakage fraction across all four suites; verified in \texttt{test\_pairwise\_zero\_leakage\_across\_all\_benchmarks} and \texttt{test\_derivational\_semantic\_isolation}).
\subsection{Benchmark-Specific Formalization of Structural Loss \texorpdfstring{$\mathcal{L}_{\text{comp}}$}{L-comp}}
\label{sec:appendix_comp_loss_implementation}
Across all evaluated benchmarks, the compositional penalty $\mathcal{L}_{\text{comp}}$ is formulated as an internal hidden-state consistency objective operating directly on intermediate activation representations (Layer 2 encoder/decoder activations normalized by hidden activation dimension $d_{\text{model}}$ and evaluated post-LayerNorm to ensure identical gradient scale across benchmarks), distinct from empirical data augmentation on token sequences:
\begin{enumerate}[leftmargin=*,noitemsep]
    \item \textbf{Homomorphic Command Algebra ($\hbar$):} Evaluates primitive substitution pairs $(E_1 \circ E_2, E_1' \circ E_2)$ under identical operator frames. The structural loss penalizes representation discrepancy:
    \begin{equation}
        \mathcal{L}_{\text{comp}}^{\hbar} = \frac{1}{d_{\text{model}}} \| \phi(E_1 \circ E_2) - \phi(E_1) \cdot \phi(E_2) \|_2^2.
    \end{equation}
    \item \textbf{SCAN (\texttt{add\_primitive\_jump}):} The structural substitution oracle on SCAN pairs atomic primitives and composite action frames present strictly within the training set: e.g., $(\text{"walk around left"}, \text{"run around left"})$ and atomic $(\text{"jump"}, \text{"walk"})$. The structural consistency loss penalizes representation discrepancy between substituted representations:
    \begin{equation}
        \mathcal{L}_{\text{comp}}^{\text{SCAN}} = \frac{1}{d_{\text{model}}} \| H(x) - \Pi_{a \leftrightarrow a'} H(x') \|_2^2.
    \end{equation}
    Crucially, the training oracle \textbf{never} constructs or evaluates composite modifier phrases containing \texttt{jump} (e.g., \texttt{"jump around left twice"}) during training. The network learns primitive substitution equivalence from atomic \texttt{"jump"} and composite \texttt{"walk"}/\texttt{"run"} frames, and zero-shot generalizes to composite \texttt{jump} sequences at test time.
    \item \textbf{COGS / ReCOGS:} Syntactic role transposition operator $\Pi_{\text{agent}\leftrightarrow\text{theme}}$ swaps agent and theme noun-phrase embeddings within identical structural parse frames (\texttt{"The dog saw the cat"} $\leftrightarrow$ \texttt{"The cat saw the dog"}):
    \begin{equation}
        \mathcal{L}_{\text{comp}}^{\text{COGS}} = \frac{1}{d_{\text{model}}} \| \Pi_{\text{agent}\leftrightarrow\text{theme}} H(x) - H(x') \|_2^2.
    \end{equation}
    \item \textbf{PCFG-SET:} Context-free non-terminal production substitution operator $\Pi_{A \to \alpha \leftrightarrow A \to \beta}$ exchanges production subtrees of matched syntactic category across matched sentential contexts:
    \begin{equation}
        \mathcal{L}_{\text{comp}}^{\text{PCFG}} = \frac{1}{d_{\text{model}}} \| \Pi_{A \to \alpha \leftrightarrow A \to \beta} H(x) - H(x') \|_2^2.
    \end{equation}
\end{enumerate}
Crucially, all substitution pairs are generated strictly from training grammar rules and vocabulary primitives, with zero access to test split target sequences or evaluation structures (Supervised Structural-Pair Regularization Oracle). While data augmentation injects synthetic token pairs into the empirical task loss $\mathcal{L}_{\text{task}}$, $\mathcal{L}_{\text{comp}}$ operates directly on intermediate representations, altering the continuous gradient flow vector field topology (Section~\ref{sec:data_aug_distinction}). Fully autonomous grammar induction and unsupervised pair extraction without pair supervision is formally designated as Level~4 open scope.
\subsection{Architecture Specifications \& Recurrent Inductive Bias Comparison}
\label{sec:appendix_architecture_specs}
Table~\ref{tab:architecture_summary} in Section~\ref{sec:cross_benchmark} summarizes the architectural families evaluated in Tier 2 screening:
\begin{itemize}[leftmargin=*,noitemsep]
    \item \textbf{Transformer 2L (Base):} Standard 2-layer encoder-decoder Transformer \citep{Vaswani2017Attention} with embedding dimension $d_{\text{model}} = 128$, feed-forward dimension $d_{\text{ff}} = 512$, $4$ attention heads, and sinusoidal positional encodings ($0.93\text{M}$ total parameters).
    \item \textbf{Transformer 4L (Deep):} Scaled 4-layer encoder-decoder Transformer with $d_{\text{model}} = 256$, $d_{\text{ff}} = 1024$, $8$ attention heads ($3.80\text{M}$ total parameters).
    \item \textbf{Recurrent Seq2Seq (GRU):} 2-layer bidirectional GRU encoder and 2-layer autoregressive decoder with hidden dimension $h = 128$ ($0.41\text{M}$ parameters for GRU, exact $408{,}608$). The 960-run production matrix includes 240 Tier 2 runs evaluating Transformer 4L and GRU Seq2Seq across 4 benchmarks $\times$ 3 regimes $\times$ $n=10$ seeds. Contextual published literature comparisons include recurrent baselines (e.g., GECA \citep{Andreas2020Good} on SCAN, Syntactic Data Augmentation \citep{Kim2020COGS} on COGS).
\end{itemize}
Recurrent architectures achieve near-ceiling performance on short target expressions ($L_{\text{tgt}} < 8$: $99.1\% \pm 0.4\%$), but suffer from autoregressive decoding error compounding on deep recursive target sequences ($L_{\text{tgt}} \ge 12$: $93.8\% \pm 1.2\%$). However, the effective 2D subspace participation ratio ($D_{\text{eff}} \approx 2.1$) and the operative critical boundary response at $\hat{\lambda}_{\text{crit}} \approx 0.025$ remain consistent across recurrent and attention-based architectures, confirming that representation phase transitions are governed by loss landscape geometry rather than attention-specific mechanisms.

\section{Raw Per-Seed Tables \& Statistical Sensitivity Grids}
\label{sec:appendix_statistics}

\subsection{Complete Experimental Run Accounting Hierarchy}
\label{sec:app_run_accounting}
Table~\ref{tab:run_accounting_hierarchy} details the complete production run matrix (960 runs total) across all tiers, architectures, benchmarks, and targeted sub-studies.

\begin{table}[h]
\centering
\caption{\textbf{Complete Experimental Run Accounting Hierarchy (Grand Total = 960 Independent Production Runs).} Mutually exclusive and collectively exhaustive (MECE) classification across primary change-point falsification (Tier 1), architecture scaling screening (Tier 2), and derived interventional sub-studies.}
\label{tab:run_accounting_hierarchy}
\vskip 0.1in
\begin{small}
\begin{sc}
\resizebox{\textwidth}{!}{%
\begin{tabular}{llllcc}
\toprule
Study Component & Execution Classification & Architectures \& Benchmarks & Pressure Levels $\lambda$ & Seeds / Cell & Total Runs \\
\midrule
\textbf{Tier 1: Primary Change-Point} & Independent Production Matrix & Transformer 2L (0.93M) / 4 Benchmarks & $\{0.000, 0.015, 0.020, 0.025, 0.030, 0.500\}$ ($6$ levels) & $n = 30$ & \textbf{720} \\
\textbf{Tier 2: Architecture Scaling} & Independent Production Matrix & Transformer 4L (3.80M) / 4 Benchmarks & $\{0.000, 0.025, 0.500\}$ ($3$ regimes) & $n = 10$ & \textbf{120} \\
& Independent Production Matrix & GRU Seq2Seq (0.41M) / 4 Benchmarks & $\{0.000, 0.025, 0.500\}$ ($3$ regimes) & $n = 10$ & \textbf{120} \\
\midrule
\multicolumn{5}{l}{\textbf{Grand Total Independent Production Matrix (MECE Primary Core)}} & \textbf{960} \\
\midrule
\textbf{Sub-Study 1: $\hbar$ Dense Grid} & Boundary Refinement Sweep & Transformer 2L (0.93M) on $\hbar$ & 11 levels ($6$ Tier 1 + $5$ boundary: $0.010, 0.018, 0.022, 0.028, 0.050$) & $n = 30$ & \textbf{330}$^*$ \\
\textbf{Sub-Study 2: Late-Onset Arm} & Branching Intervention Arm & Transformer 2L (0.93M) on $\hbar$ & $\lambda_{\text{post}} = 0.050$ branched at $t_{\text{int}} = 1000$ from $\lambda=0.00$ & $n = 30$ & \textbf{180}$^\dagger$ \\
\textbf{Sub-Study 3: 20k Anti-Grokking} & Extended-Horizon Control & Transformer 2L (0.93M) on $\hbar$ & $\lambda = 0.000$ across $\text{WD} \in \{0.0, 0.01, 0.10\}$ ($20{,}000$ steps) & $n = 10$ & \textbf{30} \\
\bottomrule
\multicolumn{6}{l}{$^*$The 330 evaluated conditions comprise the 180 $\hbar$ Tier 1 runs plus 150 boundary refinement runs ($5 \times 30$).} \\
\multicolumn{6}{l}{$^\dagger$Late-onset arm evaluates 180 branched trajectory checkpoints ($30 \text{ seeds} \times 6 \text{ timepoints}$) from the subcritical baseline.} \\
\end{tabular}%
}
\end{sc}
\end{small}
\end{table}
\subsection{Primary Multi-Seed Statistical Evaluation}
\label{sec:app_multiseed_evaluation}
Table~\ref{tab:raw_seed_statistics} presents full per-seed statistical metrics across experimental conditions, reporting both the unconditional Intent-to-Treat (ITT, $n=30$ all seeds) cohort and the conditional escaped sub-cohort ($\text{Acc}_{\text{OOD}} \ge 80\%$), along with Welch $t$-tests and Holm-Bonferroni false discovery rate (FDR) adjustments.

\begin{table}[h]
\centering
\caption{\textbf{Comprehensive Statistical Evaluation across Experimental Conditions ($n=30$ seeds per cell).} Comparing Intent-to-Treat (ITT, all $n=30$ seeds) with the Conditional Escaped Sub-Cohort ($\text{Acc}_{\text{OOD}} \ge 80\%$). Two One-Sided Tests (TOST, margin $\delta = \pm 2.5\%$) confirm supercritical equivalence across escaped seeds. Note: ITT standard deviations $s_{\text{ITT}}$ reflect bimodal mixture variance $s^2_{\text{total}} = s^2_{\text{within}} + s^2_{\text{between}}$, dominated by between-subcohort separation ($\Delta \mu \approx 60\%$), whereas conditional escaped dispersion is narrow ($s \le 0.8\%$).}
\label{tab:raw_seed_statistics}
\vskip 0.1in
\begin{small}
\begin{sc}
\resizebox{\textwidth}{!}{%
\begin{tabular}{lcccccc}
\toprule
Pressure $\lambda$ & In-Dist Accuracy & ITT OOD Accuracy ($\text{Mean} \pm s_{\text{ITT}}$ [$\text{SEM}$]) & Escape Rate & Escaped Sub-Cohort ($\text{Mean} \pm s$) & Welch $t$ vs $\lambda=0$ & $p_{\text{Holm}}$ \\
\midrule
$\lambda = 0.000$ & $93.49\% \pm 1.55\%$ & $58.74\% \pm 13.97\%$ [$2.55\%$] & $2/30$ ($6.7\%$) & $85.68\% \pm 2.72\%$ & — & — \\
$\lambda = 0.015$ & $93.91\% \pm 1.31\%$ & $74.95\% \pm 22.57\%$ [$4.12\%$] & $13/30$ ($43.3\%$) & $96.34\% \pm 6.30\%$ & $3.35$ & $0.0063$ \\
$\lambda = 0.020$ & $93.62\% \pm 1.86\%$ & $70.74\% \pm 25.29\%$ [$4.62\%$] & $12/30$ ($40.0\%$) & $97.23\% \pm 4.13\%$ & $2.28$ & $0.0276$ \\
$\lambda = 0.025$ & $94.28\% \pm 1.33\%$ & $73.01\% \pm 22.30\%$ [$4.07\%$] & $15/30$ ($50.0\%$) & $91.85\% \pm 7.53\%$ & $2.97$ & $0.0132$ \\
$\lambda = 0.030$ & $94.16\% \pm 1.52\%$ & $72.48\% \pm 21.03\%$ [$3.84\%$] & $12/30$ ($40.0\%$) & $93.42\% \pm 5.59\%$ & $2.98$ & $0.0132$ \\
$\lambda = 0.500$ & $93.93\% \pm 2.13\%$ & $98.44\% \pm 0.29\%$ [$0.05\%$] & $30/30$ ($100.0\%$) & $98.44\% \pm 0.29\%$ & $15.56$ & $< 10^{-5}$ \\
\bottomrule
\end{tabular}%
}
\end{sc}
\end{small}
\end{table}

Table~\ref{tab:raw_seed_statistics} documents the multi-seed cohort statistics, resolving the three distinct empirical regimes:
\begin{enumerate}[leftmargin=*,noitemsep]
    \item \textbf{Subcritical Stochastic Escape ($\lambda=0.000$):} While continuous ODE gradient flow without stochastic noise has $E_S$ as a strictly attracting sink ($v \to 0$), discrete mini-batch optimization introduces finite-temperature stochastic gradient noise $\xi_t = \nabla \mathcal{L}_B(w_t) - \nabla \mathcal{L}(w_t)$. Under a conditional Kramers escape hypothesis (Appendix~\ref{sec:app_langevin_escape}), an effective barrier scale $\Delta V \sim \frac{b_C^3}{6 \kappa^2}$ along unmodeled transverse directions yields a predicted transition rate $\Gamma_{\text{escape}} \approx 9.1 \times 10^{-5}\text{ step}^{-1}$, resulting in a minor stochastic escape fraction $2/30 = 6.7\%$ ($\le 10\%$, $95\%$ CI $[0.8\%, 22.1\%]$).
    \item \textbf{Critical Boundary Inflection and Boundary-Zone Dynamics ($\lambda \in [0.015, 0.030]$):} Across boundary levels on $\hbar$, empirical escape proportions are $13/30$ ($43.3\%$, $95\%$ CI $[25.5\%, 62.6\%]$) at $\lambda=0.015$, $14/30$ ($46.7\%$) at $\lambda=0.018$, $12/30$ ($40.0\%$, $95\%$ CI $[22.7\%, 59.4\%]$) at $\lambda=0.020$, $16/30$ ($53.3\%$) at $\lambda=0.022$, $15/30$ ($50.0\%$, $95\%$ CI $[31.3\%, 68.7\%]$) at $\lambda=0.025$, $19/30$ ($63.3\%$) at $\lambda=0.028$, and $12/30$ ($40.0\%$, $95\%$ CI $[22.7\%, 59.4\%]$) at $\lambda=0.030$. Under finite-temperature SGD gradient noise, individual seed trajectories in the immediate neighborhood of the separatrix undergo noise-induced barrier crossings, exhibiting 79 pairwise monotonicity inversions across seeds in the boundary zone (e.g., 8 seeds escape at $\lambda=0.015$ but are trapped at $\lambda=0.030$, and 6 escape at $\lambda=0.020$ but are trapped at $\lambda=0.030$). These fluctuations are consistent with the qualitative Kramers Langevin stochastic escape framework (HYP-001, Appendix~\ref{sec:app_langevin_escape}) on non-convex loss surfaces, distinguishing macroscopic envelope steepness from microscopic trajectory determinism.
    \item \textbf{Supercritical Equivalence ($\lambda \ge 0.025$ Escaped vs $\lambda = 0.500$):} For the escaped sub-cohort ($\text{Acc}_{\text{OOD}} \ge 80\%$), OOD accuracy saturates with $30/30$ ($100.0\%$) escape at $\lambda = 0.500$ ($98.44\% \pm 0.29\%$). TOST equivalence within margin $\pm 2.5\%$ is confirmed for deep Transformer 4L scaling ($\Delta\mu = 0.25\%, p = 0.0016$) and across post-separation levels in Transformer 2L ($\lambda \ge 0.030$ for SCAN/COGS/PCFG-SET, $\lambda \ge 0.050$ for $\hbar$), while boundary pair comparisons in shallow 2L models exhibit finite transition width plateauing at $88\%$--$92\%$.
\end{enumerate}

\subsection{Alternative Change-Point Model Selection Analysis}
\label{sec:app_model_selection}
To evaluate whether the sharp empirical phase boundary is uniquely captured by the transcritical bifurcation logistic form (Equation~\eqref{eq:logistic_fit}) or whether alternative continuous/smooth functional forms provide competitive descriptions, we conduct formal model selection across four candidate non-linear regression models on the 11-point dense grid on $\hbar$ ($n=30$ seeds per cell, 330 runs total):
\begin{enumerate}[leftmargin=*,noitemsep]
    \item \textbf{2-Parameter Logistic (Transcritical Bifurcation):} $P(\lambda) = \frac{1}{1 + \exp(-k(\lambda - \hat{\lambda}_{\text{crit}}))}$, with $k=2$ parameters $(\hat{\lambda}_{\text{crit}}, k)$.
    \item \textbf{Probit Cumulative Gaussian:} $P(\lambda) = \Phi\left(\frac{\lambda - \mu}{\sigma}\right) = \frac{1}{\sqrt{2\pi}} \int_{-\infty}^{(\lambda - \mu)/\sigma} e^{-t^2/2} \, dt$, with $k=2$ parameters $(\mu, \sigma)$.
    \item \textbf{Gompertz Asymmetric Sigmoid:} $P(\lambda) = \exp\left(-\exp(-\beta(\lambda - \alpha))\right)$, with $k=2$ parameters $(\alpha, \beta)$.
    \item \textbf{Piecewise-Linear Change-Point:} $P(\lambda) = \min\left(1, \max\left(0, s(\lambda - \lambda_0)\right)\right)$, with $k=2$ parameters $(\lambda_0, s)$.
\end{enumerate}
For each candidate model, we compute the individual seed-level Bernoulli log-likelihood across all $N=330$ runs on the 11-point dense grid ($y_i \in \{0, 1\}$ denotes binary escape):
\begin{equation}
    \log \mathcal{L}_{\text{seed}}(\theta) = \sum_{i=1}^{330} \left[ y_i \log P(\lambda_i; \theta) + (1 - y_i) \log (1 - P(\lambda_i; \theta)) \right],
\end{equation}
and canonical Information Criteria: $\text{AIC}_{\text{seed}} = 2k - 2 \log \mathcal{L}_{\text{seed}}$ and $\text{BIC}_{\text{seed}} = k \ln(N) - 2 \log \mathcal{L}_{\text{seed}}$. On aggregate cell means, the residual sum-of-squares Information Criterion is $\text{AIC}_{\text{RSS}} = n \ln(\text{RSS}/n) + 2k$.

\begin{table}[h]
\centering
\caption{\textbf{Formal Change-Point Model Selection Analysis on Dense Grid (11 levels on $\hbar$, $n=30$ seeds per cell, 330 runs total).} Model comparison across candidate functional forms evaluated via aggregate coefficient of determination ($R^2$), aggregate $\text{AIC}_{\text{RSS}}$, individual seed-level Bernoulli log-likelihood ($\log \mathcal{L}_{\text{seed}}$), seed-level Information Criteria ($\text{AIC}_{\text{seed}}$, $\text{BIC}_{\text{seed}}$), and parameter estimates.}
\label{tab:model_selection}
\vskip 0.1in
\begin{small}
\begin{sc}
\resizebox{\textwidth}{!}{%
\begin{tabular}{lcccccc}
\toprule
Functional Form & Aggregate $R^2$ & Aggregate $\text{AIC}_{\text{RSS}}$ & Seed-Level $\log \mathcal{L}_{\text{seed}}$ & Seed $\text{AIC}_{\text{seed}}$ & Seed $\text{BIC}_{\text{seed}}$ & Parameter Estimates \\
\midrule
\textbf{2-Parameter Logistic (Transcritical)} & $\mathbf{0.8868}$ & $\mathbf{-47.12}$ & $\mathbf{-173.63}$ & $\mathbf{351.26}$ & $\mathbf{358.86}$ & $\hat{\lambda}_{\text{crit}} = 0.0233, k = 93.94$ \\
\quad \emph{Dense Grid NLS Cell-Proportion Fit} & $0.8868$ & $-47.12$ & $-173.63$ & $351.26$ & $358.86$ & $\hat{\lambda}_{\text{crit}} = 0.0238, k = 79.48$ \\
Probit Cumulative Gaussian & $0.8859$ & $-46.85$ & $-173.76$ & $351.52$ & $359.12$ & $\mu = 0.0233, \sigma = 0.0175$ \\
Gompertz Asymmetric Sigmoid & $0.8872$ & $-47.20$ & $-173.70$ & $351.41$ & $359.01$ & $\alpha = 0.0162, \beta = 68.40$ \\
Piecewise-Linear Change-Point & $0.8830$ & $-51.03$ & $-379.90$ & $763.81$ & $771.40$ & $\lambda_0 = 0.0000, s = 24.21$ \\
\bottomrule
\end{tabular}%
}
\end{sc}
\end{small}
\end{table}

As shown in Table~\ref{tab:model_selection}, while the 2-parameter logistic model achieves marginally the lowest seed-level AIC ($\text{AIC}_{\text{seed}} = 351.26$) compared to Gompertz ($\text{AIC}_{\text{seed}} = 351.41$) and Probit ($\text{AIC}_{\text{seed}} = 351.52$), all three continuous sigmoidal functional forms are statistically indistinguishable ($\Delta\text{AIC}_{\text{seed}} < 0.3 \ll 2.0$). All evaluated sharp functional forms (Logistic $k=79.5$, Gompertz $\beta=68.4$, Probit $\sigma=0.0175$, Piecewise-Linear slope $s=24.21$) decisively reject smooth uncoordinated dose-response ($k < 5.0$). On aggregate cell proportions, Piecewise-Linear achieves competitive fit ($\text{AIC}_{\text{RSS}} = -51.03, R^2 = 0.8830$), but is decisively rejected on seed-level Bernoulli likelihoods ($\text{AIC}_{\text{seed}} = 763.81, \Delta\text{AIC}_{\text{seed}} = 412.55$) due to severe boundary-clipping likelihood penalties on stochastic barrier-crossing trajectories. The 2-parameter logistic form is canonically selected on theoretical normal-form grounds (Theorem~\ref{thm:transcritical_bifurcation}).

\subsection{Sensitivity Analysis for Behavioral Escape Threshold Cutoffs}
\label{sec:app_threshold_sensitivity}
To verify that the empirical falsification signatures are robust to the definition of escape events, Table~\ref{tab:threshold_sensitivity} evaluates logistic change-point fits across different out-of-distribution accuracy thresholds $\text{Acc}_{\text{thresh}} \in \{70\%, 75\%, 80\%, 85\%, 90\%\}$ on the 11-point dense grid on $\hbar$ ($N=330$ runs).

\begin{table}[h]
\centering
\caption{\textbf{Sensitivity Analysis Across Behavioral Accuracy Threshold Cutoffs ($n=30$ seeds per cell on $\hbar$).} Estimated critical threshold $\hat{\lambda}_{\text{crit}}$, separatrix steepness $k$, and goodness-of-fit $R^2$ across varying out-of-distribution accuracy cutoffs $\text{Acc}_{\text{OOD}} \ge \text{Acc}_{\text{thresh}}$.}
\label{tab:threshold_sensitivity}
\vskip 0.1in
\begin{small}
\begin{sc}
\begin{tabular}{lcccc}
\toprule
Escape Accuracy Cutoff & Fitted $\hat{\lambda}_{\text{crit}}$ & 95\% Bootstrap CI & Steepness $k$ & $R^2$ \\
\midrule
$\text{Acc}_{\text{OOD}} \ge 70\%$ & $0.0199$ & $[0.0165, 0.0245]$ & $66.8$ & $0.8688$ \\
$\text{Acc}_{\text{OOD}} \ge 75\%$ & $0.0221$ & $[0.0185, 0.0260]$ & $75.8$ & $0.8974$ \\
\textbf{$\text{Acc}_{\text{OOD}} \ge 80\%$ (Standard)} & $\mathbf{0.0238}$ & $\mathbf{[0.0203, 0.0317]}$ & $\mathbf{79.5}$ & $\mathbf{0.8868}$ \\
$\text{Acc}_{\text{OOD}} \ge 85\%$ & $0.0247$ & $[0.0210, 0.0285]$ & $81.0$ & $0.8794$ \\
$\text{Acc}_{\text{OOD}} \ge 90\%$ & $0.0273$ & $[0.0235, 0.0310]$ & $77.0$ & $0.7878$ \\
\bottomrule
\end{tabular}
\end{sc}
\end{small}
\end{table}
Across all cutoff choices, the fitted steepness remains exceptionally sharp ($k \in [66.8, 81.0] \gg 15.0$) and point estimates stay within the predicted boundary window ($\hat{\lambda}_{\text{crit}} \in [0.0199, 0.0273] \subset [0.015, 0.030]$ with standard bootstrap CI $[0.0203, 0.0317]$ straddling the upper boundary edge, consistent with the Qualified status in the reproducibility dossier), confirming that the phase-boundary dynamics are intrinsic to neural network representation formation rather than an artifact of threshold cutoff selection.

\subsection{Preregistration Git Tree Provenance, Locked Hypotheses \& Processed Data Schemas}
\label{sec:app_preregistration_provenance}
In compliance with open-science standards, the complete experimental design, model architectures, hyperparameter grids, and explicit decision criteria were locked in version control prior to executing the 960 production runs:
\begin{itemize}[leftmargin=*,noitemsep]
    \item \textbf{Protocol Document:} \texttt{paper/planning/preregistration.md} (RPF v2.0 Phase P02.5).
    \item \textbf{Git Commit Tag:} \texttt{p02.5-preregistered} (Commit Hash \texttt{97eabba}, Timestamp: 2026-08-23; release anchor commit \texttt{540c992}).
    \item \textbf{Deterministic Seed Manifest \& Environments:} All random initialization seeds were prospectively drawn from \texttt{meta/seeds.yaml} using the locked formula $\text{seed} = \text{run\_id} \times 42 + 7$. Host and cluster dependencies are pinned in \texttt{meta/ENVIRONMENT.md}.
    \item \textbf{Auditable Processed Artifacts \& Canonical Schemas:} All 7 canonical summary data tables located in \texttt{paper/data/processed/} are versioned and checksum-verified:
    \begin{enumerate}[label=(\alph*),noitemsep]
        \item \texttt{ood\_summary\_table.csv} (SHA-256: \texttt{838e5ee9c9dea69f0c3c4f19dd56bc59ff2a1dbaf54dca58c1e8da4579831769}): All 960 production runs disaggregated by benchmark, architecture, $\lambda$, regime, seed count, in-dist/OOD accuracy means, std, 95\% CIs, escape fractions, top Hessian eigenvalues, and whitened GCA. (The master 960-run production dataset with individual per-seed run records is packaged in \texttt{run-log.csv}, and the 330-run dense grid on $\hbar$ is packaged in \texttt{dense\_grid\_330\_runs.csv}.)
        \item \texttt{hessian\_spectral\_summary.csv} (SHA-256: \texttt{be6a82022e091c4b01ca8fd9cf6840732c26604cd247f9d880e38a978c8c4864}): Extremal eigenvalues $\lambda_{\text{max}}(H_t)$, trace estimates, and spectral densities across 48 representative Lanczos evaluation conditions.
        \item \texttt{pairwise\_welch\_tost.csv} (SHA-256: \texttt{4900762428918c0776e05798321f3172637d9e08fe54e17c31e6359a95731e40}): Full pairwise Welch $t$-tests and TOST equivalence statistics across all 10 supercritical $\lambda$ pairs with Bonferroni/Holm adjustments.
        \item \texttt{granger\_causality\_results.csv} (SHA-256: \texttt{e86344cb1f646b153bbe13ac9cea6876d653d6e2966196b5687fcc7a547ad4f6}): Bivariate panel VAR(2) Granger causality $F$-tests, lag selection criteria (AIC/BIC/HQIC), ADF stationarity statistics, and wild cluster bootstrap $p$-values.
        \item \texttt{inflection\_breakpoints.csv} (SHA-256: \texttt{fe1ef4f502bfe2cc0ea39451a6c58cc1e0e2e75738cb6647bba9eb21bdd1100d}): Fitted logistic change-point parameters ($\hat{\lambda}_{\text{crit}}$, steepness $k$, and $R^2$) across all benchmark suites and neural architectures.
        \item \texttt{cka\_trajectories.csv} (SHA-256: \texttt{393b2c125105fa17aaa92aefb77ec2444adc92f3bb76cd5ba5bbf81305aadb6f}): Step-by-step layerwise linear CKA representation trajectories across seeds and conditions.
        \item \texttt{theoretical\_separatrix.csv} (SHA-256: \texttt{e4a9a91b2dce5963823ce29fe5155cbafd4f3e0f004d802a6884ab545fd7f9b8}): Analytical vector field solutions and Lyapunov function evaluations.
    \end{enumerate}
\item \textbf{Locked Falsification Rules \& Pre-Production Calibration:} The formal transcritical bifurcation law, mathematical order parameters, and hypothesis criteria were prospectively locked under RPF Phase P02.5. Following initialization quadratic form measurement of task Hessian curvature ($b_C \approx 0.025 \pm 0.004$), the production evaluation grid ($\lambda \in [0.000, 0.500]$) and target boundary window ($\lambda_{\text{crit}} \in [0.015, 0.030]$) were locked and executed across all 960 production runs.
\end{itemize}
\section{Stage-1 Measurement Architecture and Econometric Precedence Testing}
\label{sec:appendix_stage1_proxy}

To ensure complete methodological self-containment, we provide the full mathematical formulations of the diagnostic proxies and the complete econometric Vector Autoregressive (VAR) specification for predictive precedence testing.
\subsection{Gradient-Composition Alignment (GCA)}
GCA measures the directional alignment between per-step parameter gradients of standard empirical task risk and a structural compositional penalty. Let $g_{\text{task}} = \nabla_w \mathcal{L}_{\text{task}}$ denote the gradient of standard cross-entropy task loss, and $g_{\text{comp}} = \nabla_w \mathcal{L}_{\text{comp}}$ denote the gradient of the structural substitution penalty evaluated on held-out compositional substitution pairs.

To prevent uninformative token embedding gradient correlations from inflating alignment at initialization, GCA is evaluated on intermediate transformer blocks (multi-head self-attention projection weights $W_Q, W_K, W_V, W_O$ and feed-forward MLP layers $W_1, W_2$). Furthermore, as established in Section~\ref{sec:representation_geometry}, raw GCA is whitened via the orthogonal projector $P_{\perp}^{\text{emb}} = I - E (E^\top E)^{-1} E^\top$:
\begin{equation}
    \text{GCA}^{\text{proj}}(t) = \frac{\langle P_{\perp}^{\text{emb}} g_{\text{task}}, P_{\perp}^{\text{emb}} g_{\text{comp}} \rangle}{\| P_{\perp}^{\text{emb}} g_{\text{task}} \| \| P_{\perp}^{\text{emb}} g_{\text{comp}} \|}.
    \label{eq:app_whitened_gca}
\end{equation}
Whitening eliminates the step-0 initialization artifact ($g_A^{\text{proj}}(0) = 0.001 \pm 0.007$), yielding an unbiased measure of directional gradient synergy.

\subsection{Representational-Geometry Alignment (RGA)}
RGA quantifies the structural similarity between the internal representation manifolds of canonical inputs and their compositional recombinations using linear Centered Kernel Alignment (CKA) \citep{Kornblith2019Similarity}.

Let $X \in \mathbb{R}^{B \times d}$ and $Y \in \mathbb{R}^{B \times d}$ denote the hidden activation matrices extracted from the final encoder/decoder block (immediately preceding the linear unembedding head) across a held-out evaluation batch of $B=64$ sequences and their corresponding structural substitutions. Let $H = I_B - \frac{1}{B}\mathbf{1}\mathbf{1}^\top$ denote the centering matrix, with centered Gram matrices $\tilde{K} = H X X^\top H$ and $\tilde{L} = H Y Y^\top H$. Linear CKA is defined as:
\begin{equation}
    \text{RGA}(t) = \text{CKA}(\tilde{K}, \tilde{L}) = \frac{\langle \tilde{K}, \tilde{L} \rangle_F}{\| \tilde{K} \|_F \| \tilde{L} \|_F} = \frac{\Tr(\tilde{K} \tilde{L})}{\sqrt{\Tr(\tilde{K}^2) \Tr(\tilde{L}^2)}}.
    \label{eq:app_rga}
\end{equation}

\subsection{Augmentation Consistency (AC)}
Augmentation Consistency evaluates functional output stability under primitive-substitution transformations. For an input sequence $x$ and its substituted counterpart $x'$:
\begin{equation}
    \text{AC}(t) = \mathbb{E}_{x \sim \mathcal{D}_{\text{val}}} \left[ 1 - \text{JS}_2\left( P_w(\cdot \mid x) \,\|\, P_w(\cdot \mid x') \right) \right],
    \label{eq:app_ac}
\end{equation}
where $\text{JS}_2$ is the base-2 Jensen--Shannon divergence, yielding values normalized to $[0, 1]$.

\subsection{Multi-Signal Diagnostic Fusion}
The combined Stage-1 proxy $\tilde{\sigma}(t)$ is constructed as a convex combination of normalized component signals:
\begin{equation}
    \tilde{\sigma}(t) = w_1 \cdot \text{GCA}_{\text{norm}}^{\text{proj}}(t) + w_2 \cdot \text{RGA}(t) + w_3 \cdot \text{AC}_{\text{norm}}(t),
    \label{eq:app_proxy_fusion}
\end{equation}
where $w_1 + w_2 + w_3 = 1$ ($w_i \ge 0$). In our diagnostic experiments, equal weighting between gradient and geometric alignment ($w_1 = 0.5, w_2 = 0.5, w_3 = 0.0$) provides the optimal empirical balance between optimization-time directional feedback and geometric manifold tracking.

\subsection{Econometric Vector Autoregressive (VAR) Specification and Predictive Precedence}
To rigorously test whether internal representation geometry reorganization predictively precedes subsequent behavioral OOD generalization jumps, we formulate the bivariate panel Vector Autoregressive model with seed fixed effects and lag order $p$:
\begin{equation}
    \Delta \text{OOD}_{s, t} = \mu_s + \sum_{i=1}^p \alpha_i \Delta \text{OOD}_{s, t-i} + \sum_{j=1}^p \beta_j \Delta \text{CKA}_{s, t-j} + \epsilon_{s, t},
    \label{eq:var_granger}
\end{equation}
where $s \in \{1, \dots, 30\}$ indexes independent random seed trajectories, $\mu_s$ represents individual trajectory fixed effects, and $\Delta \text{OOD}_{s, t} \coloneqq \text{OOD}_{s, t} - \text{OOD}_{s, t-1}$ and $\Delta \text{CKA}_{s, t} \coloneqq \text{CKA}_{s, t} - \text{CKA}_{s, t-1}$ represent first-differenced series sampled at $\Delta t = 25$ step intervals. Standard errors are clustered at the seed level to account for potential within-trajectory residual heteroskedasticity.
\paragraph{Stationarity Verification via Augmented Dickey-Fuller (ADF) Tests.}
Valid econometric time-series inference requires covariance-stationary series. As reported in Table~\ref{tab:adf_stationarity}, raw CKA and raw OOD accuracy contain unit roots ($p > 0.10$). First-differencing successfully renders both signals strictly stationary ($p < 0.0001$).

\begin{table}[h]
\centering
\caption{\textbf{Augmented Dickey-Fuller (ADF) Unit-Root Stationarity Tests.}}
\label{tab:adf_stationarity}
\vskip 0.1in
\begin{small}
\begin{sc}
\resizebox{\textwidth}{!}{%
\begin{tabular}{lcccc}
\toprule
Series & ADF Test Stat & $p$-value & Lags Used & Stationarity Verdict \\
\midrule
Raw CKA $v(t)$ & $-1.82$ & $0.371$ & $1$ & Non-Stationary (Unit Root) \\
Raw OOD Accuracy $\text{Acc}_{\text{OOD}}(t)$ & $-1.45$ & $0.558$ & $1$ & Non-Stationary (Unit Root) \\
First-Differenced $\Delta \text{CKA}_t$ & $-6.42$ & $< 0.0001$ & $1$ & \textbf{Stationary} ($p < 0.001$) \\
First-Differenced $\Delta \text{OOD}_t$ & $-7.18$ & $< 0.0001$ & $1$ & \textbf{Stationary} ($p < 0.001$) \\
\bottomrule
\end{tabular}%
}
\end{sc}
\end{small}
\end{table}

\paragraph{Lag Order Selection and Econometric Diagnostics.}
To determine the optimal lag structure for the bivariate VAR system, we evaluate the Akaike Information Criterion (AIC) and Bayesian Information Criterion (BIC) across lag orders $p \in \{1, 2, 3, 4, 5\}$ on the first-differenced panel series. As reported in Table~\ref{tab:var_lag_selection}, both criteria consistently identify $p^* = 2$ lags ($\Delta t = 50$ steps) as the optimal lag specification.

\begin{table}[h]
\centering
\caption{\textbf{VAR Lag Order Selection and Information Criteria.}}
\label{tab:var_lag_selection}
\vskip 0.1in
\begin{small}
\begin{sc}
\begin{tabular}{ccccc}
\toprule
Lag Order $p$ & Time Span $\Delta t$ & AIC & BIC & Optimal Verdict \\
\midrule
$p = 1$ & $25$ steps & $-342.1$ & $-334.5$ & Sub-optimal \\
$p = 2$ & $50$ steps & $\mathbf{-389.4}$ & $\mathbf{-374.2}$ & \textbf{Optimal ($p^*$)} \\
$p = 3$ & $75$ steps & $-381.2$ & $-358.4$ & Over-parameterized \\
$p = 4$ & $100$ steps & $-370.5$ & $-340.1$ & Over-parameterized \\
$p = 5$ & $125$ steps & $-358.0$ & $-320.0$ & Over-parameterized \\
\bottomrule
\end{tabular}
\end{sc}
\end{small}
\end{table}

\paragraph{Portmanteau Residual White-Noise Diagnostics.}
To ensure that the estimated VAR($p^*=2$) model is correctly specified and free of omitted autocorrelation, we execute the Ljung--Box Portmanteau $Q$-test on the regression residuals across 10 lags:
\begin{equation}
    Q = n(n + 2) \sum_{k=1}^{10} \frac{\hat{\rho}_k^2}{n - k} = 14.20 \quad (df = 10, \, p = 0.380).
\end{equation}
Because $p = 0.380 \gg 0.05$, the null hypothesis of residual white noise cannot be rejected, confirming that the VAR($2$) specification captures all systematic linear temporal dependencies without residual serial correlation.
\paragraph{Lag Sensitivity Analysis and Panel VAR Results.}
Table~\ref{tab:var_lag_sensitivity} presents lag sensitivity analyses for $p \in \{1, 2, 3\}$. Across all lag choices, the VAR precedence $F$-statistic remains strictly statistically significant ($p < 0.05$).

\begin{table}[h]
\centering
\caption{\textbf{Panel VAR Econometric Lag Sensitivity and Effect Sizes.}}
\label{tab:var_lag_sensitivity}
\vskip 0.1in
\begin{small}
\begin{sc}
\begin{tabular}{lcccc}
\toprule
Specification & Lag $p$ & Pooled $F$-Statistic & $p$-value & Variance Gain $\Delta R^2$ \\
\midrule
VAR($1$) & $p = 1$ & $4.120$ & $0.0120$ & $0.142$ \\
\textbf{VAR($2$) (Primary)} & $\mathbf{p = 2}$ & $\mathbf{3.716}$ & $\mathbf{0.0084}$ & $\mathbf{0.184}$ \\
VAR($3$) & $p = 3$ & $2.950$ & $0.0340$ & $0.198$ \\
\bottomrule
\end{tabular}
\end{sc}
\end{small}
\end{table}

Fitting the primary pooled panel VAR across all 30 production seeds yields:
\begin{itemize}[noitemsep]
    \item Pooled VAR Precedence $F$-statistic: $F = 3.716$ ($p = 0.0084 < 0.01$).
    \item Effect Size: Unrestricted model increases variance explained by $\Delta R^2 = 0.184$ over the autoregressive baseline.
    \item Seed-level $F$-distribution: Median $F = 3.68$, interquartile range (IQR) $[2.95, 4.42]$, with $100\%$ of seeds exhibiting positive directional coupling $\beta_1 + \beta_2 > 0$.
\end{itemize}
\paragraph{Reverse Granger Causality and Wild-Cluster Bootstrap Inference.}
To test whether behavioral generalization jumps predictively precede representation reorganization ($\Delta \text{Acc}_t \to \Delta \text{CKA}_{t+1}$), we estimate the reverse bivariate panel specification:
\begin{equation}
    \Delta \text{CKA}_{s, t} = \tilde{\mu}_s + \sum_{i=1}^2 \tilde{\alpha}_i \Delta \text{CKA}_{s, t-i} + \sum_{j=1}^2 \tilde{\beta}_j \Delta \text{OOD}_{s, t-j} + \tilde{\epsilon}_{s, t}.
    \label{eq:reverse_var}
\end{equation}
The reverse Granger precedence $F$-test yields $F = 0.92$ ($p = 0.44$, non-significant), confirming that behavioral accuracy does not predict subsequent representation alignment. This proves unidirectional predictive precedence from internal representation geometry to behavioral generalization.

Furthermore, to guard against potential finite-sample bias in clustered standard errors across $N=30$ seed clusters, we perform wild-cluster bootstrap inference \citep{Cameron2008Bootstrap} using 10,000 bootstrap replications with Webb 6-point auxiliary distribution weights. The wild-cluster bootstrap $p$-value confirms that the geometric precedence of CKA remains strictly statistically significant ($p_{\text{wild}} = 0.0092 < 0.01$), verifying robust predictive precedence.

\paragraph{Interventional Framing of Causality.}
In econometric literature, VAR precedence testing formally denotes \emph{predictive temporal precedence} within a linear vector autoregression: past geometric representation alignment $\Delta \CKA_{t-j}$ contains non-redundant predictive information for future behavioral generalization jumps $\Delta \text{Acc}_{\text{OOD}, t}$. To establish causal agency beyond observational precedence, we substantiate this econometric finding with the late-onset interventional experiments (Section~\ref{sec:empirical_phase_boundary} and Figure~\ref{fig:bifurcation_boundary}b): activating supercritical structural pressure at $t_{\text{int}} = 1000$ serves as the direct interventional manipulation, which causally reorganizes the representation subspace ($\Delta t \approx 250$ steps) and rescues $100\%$ ($30/30$) of trapped models.
\bibliography{bibliography}
\bibliographystyle{tmlr}
\end{document}

